\documentclass[../libro.tex]{subfiles}

\begin{document}

\ifSubfilesClassLoaded{\appendix%
\chapter{Afterstories}\clearpage}{}

\section{Sylvester 阵与结式}

本节, 我们用行列式研究整式的一些性质.

为方便, 我们作一个有名的 (有名字的) 假定.
我们会反复地说它.

\begin{restatable}{assumption}{AssumptionResultant}%
    \label{standing-assumption-of-resultant}
    设 \(n\), \(m\) 是非负整数,
    且 \(m + n \geq 1\).
    设
    \begin{align*}
        f(x) &
        = \sum_{k = 0}^{n} {a_k x^{n-k}}
        = a_0 x^n + a_1 x^{n-1} + \dots + a_n,
        \\
        g(x) &
        = \sum_{k = 0}^{m} {b_k x^{m-k}}
        = b_0 x^m + b_1 x^{m-1} + \dots + b_m.
    \end{align*}
    为方便, 我们约定:
    若 \(k < 0\) 或 \(k > n\), 则 \(a_k = 0\);
    若 \(k < 0\) 或 \(k > m\), 则 \(b_k = 0\).
\end{restatable}

\begin{definition}
    设假定~\ref{standing-assumption-of-resultant} 成立.
    定义 \(f\), \(g\) 的 \emph{Sylvester 阵}%
    为 \(m + n\)~级阵
    \(S(f, g; n, m)\),
    其中,
    \begin{align*}
        [S(f, g; n, m)]_{i,j} =
        \begin{cases}
            a_{i-j},     & j \leq m; \\
            b_{i-(j-m)}, & j > m.
        \end{cases}
    \end{align*}
    \(f\), \(g\) 的 Sylvester 阵的行列式
    \(\det {(S(f, g; n, m))}\)
    是 \(f\), \(g\) 的\emph{结式}.
\end{definition}

我们看几个例.

\begin{example}
    设假定~\ref{standing-assumption-of-resultant} 成立.
    设 \(n = 2\), 且 \(m = 5\).
    则 \(f\), \(g\) 的 Sylvester 阵
    \(S(f, g; 2, 5)\) 是
    \begin{equation*}
        \begin{bmatrix}
            a_{0} & a_{-1} & a_{-2} & a_{-3} & a_{-4} & b_{0} & b_{-1} \\
            a_{1} & a_{0}  & a_{-1} & a_{-2} & a_{-3} & b_{1} & b_{0}  \\
            a_{2} & a_{1}  & a_{0}  & a_{-1} & a_{-2} & b_{2} & b_{1}  \\
            a_{3} & a_{2}  & a_{1}  & a_{0}  & a_{-1} & b_{3} & b_{2}  \\
            a_{4} & a_{3}  & a_{2}  & a_{1}  & a_{0}  & b_{4} & b_{3}  \\
            a_{5} & a_{4}  & a_{3}  & a_{2}  & a_{1}  & b_{5} & b_{4}  \\
            a_{6} & a_{5}  & a_{4}  & a_{3}  & a_{2}  & b_{6} & b_{5}  \\
        \end{bmatrix}
        =
        \begin{bmatrix}
            a_0 & 0   & 0   & 0   & 0   & b_0 & 0   \\
            a_1 & a_0 & 0   & 0   & 0   & b_1 & b_0 \\
            a_2 & a_1 & a_0 & 0   & 0   & b_2 & b_1 \\
            0   & a_2 & a_1 & a_0 & 0   & b_3 & b_2 \\
            0   & 0   & a_2 & a_1 & a_0 & b_4 & b_3 \\
            0   & 0   & 0   & a_2 & a_1 & b_5 & b_4 \\
            0   & 0   & 0   & 0   & a_2 & 0   & b_5 \\
        \end{bmatrix}.
        \qefhere
    \end{equation*}
\end{example}

\begin{example}
    设假定~\ref{standing-assumption-of-resultant} 成立.
    设 \(m = n = 1\).
    则 \(f(x) = a_0 x + a_1\),
    且 \(g(x) = b_0 x + b_1\).
    则 \(f\), \(g\) 的 Sylvester 阵
    \begin{align*}
        S(f, g; 1, 1) =
        \begin{bmatrix}
            a_0 & b_0 \\
            a_1 & b_1 \\
        \end{bmatrix},
    \end{align*}
    且
    \begin{equation*}
        \det {(S(f, g; 1, 1))} = a_0 b_1 - b_0 a_1.
        \qefhere
    \end{equation*}
\end{example}

\begin{example}\label{emp:morniszero}
    设假定~\ref{standing-assumption-of-resultant} 成立.
    设 \(n = 0\).
    则 \(f(x) = a_0\).
    则 \(m\)~级阵 \(S(f, g; \allowbreak
    0, m)\) 适合,
    \begin{align*}
        [S(f, g; 0, m)]_{i,j} = a_{i-j} =
        \begin{cases}
            a_0, & i = j;       \\
            0,   & \text{别的情形}.
        \end{cases}
    \end{align*}
    则 \(S(f, g; 0, m) = a_0 I_m\),
    且 \(\det {(S(f, g; 0, m))} = a_0^m\).

    类似地, 若 \(m = 0\),
    则 \(n\)~级阵 \(S(f, g; n, 0) = b_0 I_n\),
    且 \(\det {(S(f, g; n, 0))} = b_0^n\).
\end{example}

\begin{example}
    设假定~\ref{standing-assumption-of-resultant} 成立.
    设 \(a_0 = 0 = b_0\).
    则 \(k \leq 0\) 时, \(a_k = 0 = b_k\).
    注意, 当 \(j \leq m\) 时,
    \([S(f, g; n, m)]_{1,j} \allowbreak
    = a_{1-j}\),
    而 \(1 - j \leq 0\),
    则 \([S(f, g; n, m)]_{1,j} \allowbreak
    = 0\).
    当 \(j > m\) 时,
    \([S(f, g; n, m)]_{1,j} \allowbreak
    = b_{1 - (j - m)}\),
    而 \(1 - (j - m) \leq 0\),
    则 \([S(f, g; n, m)]_{1,j} \allowbreak
    = 0\).
    则 \(S(f, g; n, m)\) 的行~\(1\) 的元全是 \(0\).
    则 \(\det {(S(f, g; n, m))} = 0\).
\end{example}

我们给一个研究 Sylvester 阵的可能的原因.
设假定~\ref{standing-assumption-of-resultant} 成立.
设数 \(c\) 适合, \(f(c) = 0 = g(c)\).
则
\begin{align*}
    a_0 c^n + a_1 c^{n-1} + \dots + a_n & = 0, \\
    b_0 c^m + b_1 c^{m-1} + \dots + b_m & = 0.
\end{align*}
若我们分别以 \(c^{m-1}\), \(c^{m-2}\), \(\dots\), \(1\) 乘等式~1,
且分别以 \(c^{n-1}\), \(c^{n-2}\), \(\dots\), \(1\) 乘等式~2,
我们得
\begin{align*}
    a_0 c^{m+n-1} + a_1 c^{m+n-2} + \dots + a_{n-1} c^{m} + a_n c^{m-1}   & = 0,   \\
    a_0 c^{m+n-2} + a_1 c^{m+n-3} + \dots + a_{n-1} c^{m-1} + a_n c^{m-2} & = 0,   \\
    & \dots, \\
    a_0 c^{n} + a_1 c^{n-1} + \dots + a_{n-1} c + a_n 1                   & = 0,   \\
    b_0 c^{m+n-1} + b_1 c^{m+n-2} + \dots + b_{m-1} c^{n} + b_m c^{n-1}   & = 0,   \\
    b_0 c^{m+n-2} + b_1 c^{m+n-3} + \dots + b_{m-1} c^{n-1} + b_m c^{n-2} & = 0,   \\
    & \dots, \\
    b_0 c^{m} + b_1 c^{m-1} + \dots + b_{m-1} c + b_m 1                   & = 0.
\end{align*}
回想, 当 \(k < 0\) 或 \(k > n\) 时, \(a_k = 0\),
且当 \(k < 0\) 或 \(k > m\) 时, \(b_k = 0\).
于是, 我们可写
\begin{align*}
    a_{0} c^{m+n-1} + a_{1} c^{m+n-2} + \dots + a_{m+n-2} c + a_{m+n-1} 1  & = 0,   \\
    a_{-1} c^{m+n-1} + a_{0} c^{m+n-2} + \dots + a_{m+n-3} c + a_{m+n-2} 1 & = 0,   \\
    & \dots, \\
    a_{1-m} c^{m+n-1} + a_{2-m} c^{m+n-2} + \dots + a_{n-1} c + a_{n} 1    & = 0,   \\
    b_{0} c^{m+n-1} + b_{1} c^{m+n-2} + \dots + b_{m+n-2} c + b_{m+n-1} 1  & = 0,   \\
    b_{-1} c^{m+n-1} + b_{0} c^{m+n-2} + \dots + b_{m+n-3} c + b_{m+n-2} 1 & = 0,   \\
    & \dots, \\
    b_{1-n} c^{m+n-1} + b_{2-n} c^{m+n-2} + \dots + b_{m-1} c + b_{m} 1    & = 0.
\end{align*}
作 \(1 \times (m + n)\)~阵
\(C = [c^{m+n-1}, c^{m+n-2}, \dots, c, 1]\).
则
\begin{align*}
    S(f, g; n, m)^{\mathrm{T}}\, C^{\mathrm{T}} = 0.
\end{align*}
注意, 线性方程组
\(S(f, g; n, m)^{\mathrm{T}}\, X = 0\)
有非零的解 \(X = C^{\mathrm{T}}\).
故, 由 Cramer 公式 (见第~1~章, 节~\sekcio{21})
与反证法,
\(S(f, g; n, m)^{\mathrm{T}}\) 的行列式一定是 \(0\).
故 \(\det {(S(f, g; n, m))} = 0\).

\begin{theorem}\label{thm:necessary-condition-for-the-existence-of-a-common-root}
    设假定~\ref{standing-assumption-of-resultant} 成立.
    设存在数 \(c\), 使 \(f(c) = 0 = g(c)\).
    则
    \begin{align*}
        \det {(S(f, g; n, m))} = 0.
    \end{align*}
\end{theorem}

我们进一步地研究 Sylvester 阵.

\begin{theorem}
    设假定~\ref{standing-assumption-of-resultant} 成立.
    作 \(1 \times (m+n)\)~阵 \(C_{m+n}\),
    其中, \([C_{m+n}]_{1,j} = x^{m+n-j}\).
    设 \(Y\) 是 \((m+n) \times 1\)~阵.
    记
    \begin{align*}
        u(x) = \sum_{k = 0}^{m-1} {[Y]_{k+1,1} x^{m-1-k}},
        \quad
        v(x) = \sum_{k = 0}^{n-1} {[Y]_{m+1+k,1} x^{n-1-k}}.
    \end{align*}

    (1)
    \(C_{m+n}\, S(f, g; n, m)\) 是 \(1 \times (m + n)\)~阵,
    且
    \begin{align*}
        [C_{m+n}\, S(f, g; n, m)]_{1,j} =
        \begin{cases}
            f(x)\, x^{m-j},     & j \leq m; \\
            g(x)\, x^{n-(j-m)}, & j > m.
        \end{cases}
    \end{align*}

    (2)
    \(C_{m+n}\, S(f, g; n, m)\, Y\) 是 \(1\)~级阵, 且
    \begin{align*}
        [C_{m+n}\, S(f, g; n, m)\, Y]_{1,1} = f(x)\, u(x) + g(x)\, v(x).
    \end{align*}

    (3)
    记
    \begin{align*}
        f(x)\, u(x) + g(x)\, v(x) = \sum_{k = 0}^{m+n-1} {c_k x^{m+n-1-k}}.
    \end{align*}
    则 \(S(f, g; n, m)\, Y\) 是 \((m+n) \times 1\)~阵, 且
    \begin{align*}
        c_k = [S(f, g; n, m)\, Y]_{k+1,1}.
    \end{align*}
\end{theorem}

\begin{proof}
    (1)
    注意,
    \begin{align*}
        [C_{m+n}\, S(f, g; n, m)]_{1,j}
        = {} &
        \sum_{\ell = 1}^{m+n}
        {[C_{m+n}]_{1,\ell} [S(f, g; n, m)]_{\ell,j}}
        \\
        = {} &
        \sum_{1 \leq \ell \leq m+n}
        {[S(f, g; n, m)]_{\ell,j} x^{m+n-\ell}}.
    \end{align*}
    若 \(j \leq m\), 则
    \([S(f, g; n, m)]_{\ell,j} = a_{\ell-j}\).
    则
    \begin{align*}
        &
        [C_{m+n}\, S(f, g; n, m)]_{1,j}
        \\
        = {} &
        \sum_{1 \leq \ell \leq m+n}
        {a_{\ell-j} x^{m+n-\ell}}
        \\
        = {} &
        \sum_{1-j \leq \ell-j \leq m+n-j}
        {a_{\ell-j} x^{n-(\ell-j)+m-j}}
        \\
        = {} &
        \sum_{1-j \leq k \leq m+n-j}
        {a_k x^{n-k+m-j}}
        \\
        = {} &
        \sum_{\substack{
                1-j \leq k \leq m+n-j \\
                0 \leq k \leq n
        }}
        {a_k x^{n-k+m-j}}
        +
        \sum_{\substack{
                1-j \leq k \leq m+n-j \\
                k < 0 \,\text{或}\, k > n
        }}
        {a_k x^{n-k+m-j}}
        \\
        = {} &
        \sum_{0 \leq k \leq n}
        {a_k x^{n-k} x^{m-j}}
        +
        \sum_{\substack{
                1-j \leq k < 0 \\
                \text{或}\, n < k \leq m+n-j
        }}
        {0\, x^{n-k+m-j}}
        \\
        = {} &
        f(x)\, x^{m-j}.
    \end{align*}
    若 \(j > m\), 则
    \([S(f, g; n, m)]_{\ell,j} = b_{\ell-(j-m)}\).
    则
    \begin{align*}
        &
        [C_{m+n}\, S(f, g; n, m)]_{1,j}
        \\
        = {} &
        \sum_{1 \leq \ell \leq m+n}
        {b_{\ell-(j-m)} x^{m+n-\ell}}
        \\
        = {} &
        \sum_{1-(j-m) \leq \ell-(j-m) \leq m+n-(j-m)}
        {b_{\ell-(j-m)} x^{m-(\ell-(j-m))+n-(j-m)}}
        \\
        = {} &
        \sum_{1-(j-m) \leq k \leq m+n-(j-m)}
        {b_k x^{m-k+n-(j-m)}}
        \\
        = {} &
        \sum_{\substack{
                1-(j-m) \leq k \leq m+n-(j-m) \\
                0 \leq k \leq m
        }}
        {b_k x^{m-k+n-(j-m)}}
        +
        \sum_{\substack{
                1-(j-m) \leq k \leq m+n-(j-m) \\
                k < 0 \,\text{或}\, k > m
        }}
        {b_k x^{m-k+n-(j-m)}}
        \\
        = {} &
        \sum_{0 \leq k \leq m}
        {b_k x^{m-k} x^{n-(j-m)}}
        +
        \sum_{\substack{
                1-(j-m) \leq k < 0 \\
                \text{或}\, m < k \leq m+n-(j-m)
        }}
        {0\, x^{m-k+n-(j-m)}}
        \\
        = {} &
        g(x)\, x^{n-(j-m)}.
    \end{align*}

    (2)
    注意,
    \begin{align*}
        &
        [C_{m+n}\, S(f, g; n, m)\, Y]_{1,1}
        \\
        = {} &
        [(C_{m+n}\, S(f, g; n, m))\, Y]_{1,1}
        \\
        = {} &
        \sum_{j = 1}^{m+n}
        {[C_{m+n}\, S(f, g; n, m)]_{1,j} [Y]_{j,1}}
        \\
        = {} &
        \sum_{j = 1}^{m}
        {[C_{m+n}\, S(f, g; n, m)]_{1,j} [Y]_{j,1}}
        +
        \sum_{j = m+1}^{m+n}
        {[C_{m+n}\, S(f, g; n, m)]_{1,j} [Y]_{j,1}}
        \\
        = {} &
        \sum_{j = 1}^{m}
        {f(x)\, x^{m-j}\, [Y]_{j,1}}
        +
        \sum_{j = m+1}^{m+n}
        {g(x)\, x^{n-(j-m)}\, [Y]_{j,1}}
        \\
        = {} &
        f(x)
        \sum_{j = 1}^{m}
        {x^{m-j}\, [Y]_{j,1}}
        +
        g(x)
        \sum_{j = m+1}^{m+n}
        {x^{n-(j-m)}\, [Y]_{j,1}}
        \\
        = {} &
        f(x)
        \sum_{j = 1}^{m}
        {x^{m-1-(j-1)}\, [Y]_{(j-1)+1,1}}
        +
        g(x)
        \sum_{j = m+1}^{m+n}
        {x^{n-1-(j-m-1)}\, [Y]_{m+1+(j-m-1),1}}
        \\
        = {} &
        f(x)
        \sum_{k = 0}^{m-1}
        {[Y]_{k+1,1} x^{m-1-k}}
        +
        g(x)
        \sum_{k = 0}^{n-1}
        {[Y]_{m+1+k,1} x^{n-1-k}}
        \\
        = {} &
        f(x)\, u(x) + g(x)\, v(x).
    \end{align*}

    (3)
    注意,
    \begin{align*}
        [C_{m+n}\, S(f, g; n, m)\, Y]_{1,1}
        = {} &
        [C_{m+n}\, (S(f, g; n, m)\, Y)]_{1,1}
        \\
        = {} &
        \sum_{j = 1}^{m+n}
        {[C_{m+n}]_{1,j} [S(f, g; n, m)\, Y]_{j,1}}
        \\
        = {} &
        \sum_{j = 1}^{m+n}
        {x^{m+n-j} [S(f, g; n, m)\, Y]_{j,1}}
        \\
        = {} &
        \sum_{j = 1}^{m+n}
        {[S(f, g; n, m)\, Y]_{(j-1)+1,1} x^{m+n-1-(j-1)}}
        \\
        = {} &
        \sum_{k = 0}^{m+n-1}
        {[S(f, g; n, m)\, Y]_{k+1,1} x^{m+n-1-k}}.
    \end{align*}
    故, 对不超过 \(m+n-1\) 的非负整数 \(k\),
    \begin{equation*}
        [S(f, g; n, m)\, Y]_{k+1,1} = c_k.
        \qedhere
    \end{equation*}
\end{proof}

\begin{theorem}
    设 \(A\) 是 \(n\)~级阵.
    则存在 \(n \times 1\)~阵 \(P\),
    使
    \begin{align*}
        [A P]_{i,1} =
        \begin{cases}
            \det {(A)}, & i = n;       \\
            0,          & \text{别的情形}.
        \end{cases}
    \end{align*}
    进一步地, 我们可要求, \(P \neq 0\).
\end{theorem}

\begin{proof}
    设 \(W\) 是 \(A\) 的古伴 \(\operatorname{adj} {(A)}\) 的列~\(n\).
    则
    \begin{align*}
        [A W]_{i,1}
        = {} &
        \sum_{j = 1}^{n} {[A]_{i,j} [W]_{j,1}}
        \\
        = {} &
        \sum_{j = 1}^{n} {[A]_{i,j} [\operatorname{adj} {(A)}]_{j,n}}
        \\
        = {} &
        [A \operatorname{adj} {(A)}]_{i,n}
        \\
        = {} &
        [\det {(A)}\, I_n]_{i,n}
        \\
        = {} &
        \det {(A)}\, [I_n]_{i,n}.
    \end{align*}
    取 \(P\) 为 \(W\) 即可.

    设, 进一步地, 我们要求, \(P \neq 0\).
    若 \(W \neq 0\), 则我们仍取 \(P\) 为 \(W\).
    若 \(W = 0\), 则
    \(\det {(A)} = [A W]_{i,n} = [A\, 0]_{i,n} = [0]_{i,n} = 0\).
    则, 由第~1~章, 节~\malneprasekcio{23},
    或由第~1~章, 节~\malneprasekcio{25},
    或由 ``REF 阵、线性方程组、秩'',
    存在非零的 \(n \times 1\) 阵 \(Q\),
    使 \(A Q = 0\).
    取 \(P\) 为 \(Q\) 即可.
\end{proof}

\begin{theorem}\label{thm:resultant-Bézout}
    设假定~\ref{standing-assumption-of-resultant} 成立.

    (1)
    存在整式 \(u\), \(v\), 使
    \begin{align*}
        f(x)\, u(x) + g(x)\, v(x) = \det {(S(f, g; n, m))}.
    \end{align*}
    注意, 上式的右边不含 \(x\)
    (因为 \(S(f, g; n, m)\) 不含 \(x\)).

    (2)
    设存在数 \(c\), 使 \(f(c) = 0 = g(c)\).
    则
    \begin{align*}
        \det {(S(f, g; n, m))} = 0.
    \end{align*}
\end{theorem}

\begin{proof}
    (1)
    我们知道, 存在 \((m + n) \times 1\)~阵 \(P\),
    使
    \begin{align*}
        [S(f, g; n, m)\, P]_{i,1} =
        \begin{cases}
            \det {(S(f, g; n, m))}, & i = m + n;   \\
            0,                      & \text{别的情形}.
        \end{cases}
    \end{align*}
    记
    \begin{align*}
        u(x) = \sum_{k = 0}^{m-1} {[P]_{k+1,1} x^{m-1-k}},
        \quad
        v(x) = \sum_{k = 0}^{n-1} {[P]_{m+1+k,1} x^{n-1-k}}.
    \end{align*}
    则
    \begin{align*}
        f(x)\, u(x) + g(x)\, v(x) = \sum_{k = 0}^{m+n-1} {c_k x^{m+n-1-k}},
    \end{align*}
    其中,
    \begin{align*}
        c_k = [S(f, g; n, m)\, P]_{k+1,1} =
        \begin{cases}
            \det {(S(f, g; n, m))}, & k = m + n - 1; \\
            0,                      & \text{别的情形}.
        \end{cases}
    \end{align*}

    (2)
    我们其实已证它
    (定理~\ref{thm:necessary-condition-for-the-existence-of-a-common-root}),
    但我们现在别地证它.
    注意, 代 \(x\) 以 \(c\), 得
    \begin{equation*}
        \det {(S(f, g; n, m))}
        = f(c)\, u(c) + g(c)\, v(c)
        = 0.
        \qedhere
    \end{equation*}
\end{proof}

有一件事, 其值得提.
设假定~\ref{standing-assumption-of-resultant} 成立.
对 \(m + n\)~级阵 \(S(f, g; n, m)\),
我们加行~\(i\) 的 \(x^{m+n-i}\)~倍于行~\(m + n\),
\(i = 1\), \(\dots\), \(m + n - 1\),
得阵 \(S^* (f, g; n, m)\).
那么, 一方面,
\(\det {(S^* (f, g; n, m))} = \det {(S(f, g; n, m))}\).
另一方面, 若 \(j \leq m\), 则
\begin{align*}
    &
    [S^* (f, g; n, m)]_{m+n,j}
    \\
    = {} &
    [S(f, g; n, m)]_{m+n,j}
    + \sum_{1 \leq \ell \leq m+n-1}
    {x^{m+n-\ell} [S(f, g; n, m)]_{\ell,j}}
    \\
    = {} &
    a_{m+n-j} x^0
    + \sum_{1 \leq \ell \leq m+n-1}
    {a_{\ell-j} x^{m+n-\ell}}
    \\
    = {} &
    \sum_{1 \leq \ell \leq m+n}
    {a_{\ell-j} x^{m+n-\ell}}
    \\
    = {} &
    \sum_{1-j \leq \ell-j \leq m+n-j}
    {a_{\ell-j} x^{n-(\ell-j)+m-j}}
    \\
    = {} &
    \sum_{1-j \leq k \leq m+n-j}
    {a_k x^{n-k+m-j}}
    \\
    = {} &
    \sum_{\substack{
            1-j \leq k \leq m+n-j \\
            0 \leq k \leq n
    }}
    {a_k x^{n-k+m-j}}
    +
    \sum_{\substack{
            1-j \leq k \leq m+n-j \\
            k < 0 \,\text{或}\, k > n
    }}
    {a_k x^{n-k+m-j}}
    \\
    = {} &
    \sum_{0 \leq k \leq n}
    {a_k x^{n-k} x^{m-j}}
    +
    \sum_{\substack{
            1-j \leq k < 0 \\
            \text{或}\, n < k \leq m+n-j
    }}
    {0\, x^{n-k+m-j}}
    \\
    = {} &
    f(x)\, x^{m-j};
\end{align*}
若 \(j > m\), 则
\begin{align*}
    &
    [S^* (f, g; n, m)]_{m+n,j}
    \\
    = {} &
    [S(f, g; n, m)]_{m+n,j}
    + \sum_{1 \leq \ell \leq m+n-1}
    {x^{m+n-\ell} [S(f, g; n, m)]_{\ell,j}}
    \\
    = {} &
    b_{m+n-(j-m)} x^0
    + \sum_{1 \leq \ell \leq m+n-1}
    {b_{\ell-(j-m)} x^{m+n-\ell}}
    \\
    = {} &
    \sum_{1 \leq \ell \leq m+n}
    {b_{\ell-(j-m)} x^{m+n-\ell}}
    \\
    = {} &
    \sum_{1-(j-m) \leq \ell-(j-m) \leq m+n-(j-m)}
    {b_{\ell-(j-m)} x^{m-(\ell-(j-m))+n-(j-m)}}
    \\
    = {} &
    \sum_{1-(j-m) \leq k \leq m+n-(j-m)}
    {b_k x^{m-k+n-(j-m)}}
    \\
    = {} &
    \sum_{\substack{
            1-(j-m) \leq k \leq m+n-(j-m) \\
            0 \leq k \leq m
    }}
    {b_k x^{m-k+n-(j-m)}}
    +
    \sum_{\substack{
            1-(j-m) \leq k \leq m+n-(j-m) \\
            k < 0 \,\text{或}\, k > m
    }}
    {b_k x^{m-k+n-(j-m)}}
    \\
    = {} &
    \sum_{0 \leq k \leq m}
    {b_k x^{m-k} x^{n-(j-m)}}
    +
    \sum_{\substack{
            1-(j-m) \leq k < 0 \\
            \text{或}\, m < k \leq m+n-(j-m)
    }}
    {0\, x^{m-k+n-(j-m)}}
    \\
    = {} &
    g(x)\, x^{n-(j-m)}.
\end{align*}
所以, 若存在数 \(c\), 使 \(f(c) = 0 = g(c)\),
则 \(S^* (f, g; n, m)\) 的行~\(m + n\) 的元全是 \(0\).
则
\begin{align*}
    \det {(S(f, g; n, m))}
    = \det {(S^* (f, g; n, m))}
    = 0.
\end{align*}

\begin{theorem}\label{thm:resultant-swap}
    设假定~\ref{standing-assumption-of-resultant} 成立.

    (1)
    \(S(g, f; m, n)\) 是 \(m + n\)~级阵,
    且 \(S(f, g; n, m)\) 的列~\(j\) 是
    \(S(g, f; m, n)\) 的列~\(n + j\)
    (若 \(j \leq m\))
    或 \(S(g, f; m, n)\) 的列~\(j - m\)
    (若 \(j > m\)).

    (2)
    \(\det {(S(g, f; m, n))}
    = (-1)^{m n} \det {(S(f, g; n, m))}\).
\end{theorem}

\begin{proof}
    (1)
    注意, \(S(g, f; m, n)\) 是这样的 \(m + n\)~级阵:
    \begin{align*}
        [S(g, f; m, n)]_{i,j} =
        \begin{cases}
            b_{i-j},     & j \leq n; \\
            a_{i-(j-n)}, & j > n.
        \end{cases}
    \end{align*}
    故, 若 \(j \leq m\),
    \begin{align*}
        [S(f, g; n, m)]_{i,j}
        = a_{i-j}
        = a_{i-(n+j-n)}
        = [S(g, f; m, n)]_{i,n+j},
    \end{align*}
    且若 \(j > m\),
    \begin{align*}
        [S(f, g; n, m)]_{i,j}
        = b_{i-(j-m)}
        = [S(g, f; m, n)]_{i,j-m}.
    \end{align*}

    (2)
    若 \(m = 0\) 或 \(n = 0\),
    则这是平凡地对的
    (见例~\ref{emp:morniszero}).

    设 \(m \geq 1\) 且 \(n \geq 1\).
    对 \(S(g, f; m, n)\),
    交换列~\(n + 1\) 与列~\(n\),
    再交换新的阵的列~\(n\) 与列~\(n - 1\),
    ……
    再交换新的阵的列~\(2\) 与列~\(1\),
    我们作了 \(n\)~次相邻的列的交换,
    使 \(S(g, f; m, n)\) 的列~\(n + 1\) 到列~\(1\),
    且别的列的相对位置不变.
    类似地, 我们也可使这个阵的%
    列~\(n + 2\), \(\dots\), \(n + m\)
    分别地到%
    列~\(2\), \(\dots\), \(m\).
    则 \(S(g, f; m, n)\) 的列 \(1\), \(\dots\), \(n\)
    分别地到%
    列~\(m + 1\), \(\dots\), \(m + n\).
    则 \(n \cdot m\)~次相邻的列的交换%
    变 \(S(g, f; m, n)\) 为 \(S(f, g; n, m)\).
    则
    \(\det {(S(g, f; m, n))}
    = (-1)^{m n} \det {(S(f, g; n, m))}\).
\end{proof}

\begin{theorem}\label{thm:resultant-recursive}
    设假定~\ref{standing-assumption-of-resultant} 成立.
    设 \(t\) 是数, 且
    \begin{align*}
        h(x)
        = f(x)\, (x - t)
        = \sum_{k = 0}^{n+1} {c_k x^{n+1-k}}.
    \end{align*}
    我们约定, 若 \(k < 0\) 或 \(k > n+1\), 则 \(c_k = 0\).
    则
    \begin{align*}
        \det {(S(h, g; n+1, m))} = \det {(S(f, g; n, m))}\, g(t),
    \end{align*}
    且
    \begin{align*}
        \det {(S(g, h; m, n+1))} = (-1)^m \det {(S(g, f; m, n))}\, g(t).
    \end{align*}
\end{theorem}

\begin{proof}
    我们证式~1;
    您用式~1 与定理~\ref{thm:resultant-swap} 证式~2.

    注意, 由整式的乘法,
    \begin{align*}
        c_k =
        \begin{cases}
            a_0,             & k = 0;           \\
            a_k - t a_{k-1}, & 1 \leq k \leq n; \\
            -t a_n,          & k = n + 1.
        \end{cases}
    \end{align*}
    回想, 若 \(k < 0\) 或 \(k > n\),
    则 \(a_k = 0\).
    于是, 统一地, \(c_k = a_k - t a_{k-1}\).
    注意, \(a_k = c_k + t a_{k-1}\).

    我们记
    \(g_\ell (x) = 0\)
    (若 \(\ell < 0\)),
    且
    \(g_\ell (x) = g_{\ell-1} (x)\, x + b_\ell\)
    (若 \(\ell \geq 0\)).
    则对任何整数 \(\ell\),
    \(g_\ell (t) = b_\ell + t g_{\ell-1} (t)\).
    注意, \(b_\ell = g_\ell (t) - t g_{\ell-1} (t)\).

    注意, \(m + n + 1\) 级阵 \(S(h, g; n+1, m)\) 等于
    \begin{align*}
        \begin{bmatrix}
            c_{0}   & c_{-1}    & c_{-2}    & \cdots & c_{1-m}   & b_{0}   & b_{-1}    & \cdots & b_{-n}  \\
            c_{1}   & c_{0}     & c_{-1}    & \cdots & c_{2-m}   & b_{1}   & b_{0}     & \cdots & b_{1-n} \\
            c_{2}   & c_{1}     & c_{0}     & \cdots & c_{3-m}   & b_{2}   & b_{1}     & \cdots & b_{2-n} \\
            \vdots  & \vdots    & \vdots    & {}     & \vdots    & \vdots  & \vdots    & {}     & \vdots  \\
            c_{n}   & c_{n-1}   & c_{n-2}   & \cdots & c_{n+1-m} & b_{n}   & b_{n-1}   & \cdots & b_{0}   \\
            c_{n+1} & c_{n}     & c_{n-1}   & \cdots & c_{n+2-m} & b_{n+1} & b_{n}     & \cdots & b_{1}   \\
            \vdots  & \vdots    & \vdots    & {}     & \vdots    & \vdots  & \vdots    & {}     & \vdots  \\
            c_{m+n} & c_{m+n-1} & c_{m+n-2} & \cdots & c_{n+1}   & b_{m+n} & b_{m+n-1} & \cdots & b_{m}   \\
        \end{bmatrix},
    \end{align*}
    即
    \begin{align*}
        \begin{bmatrix}
            a_{0}   & a_{-1}    & a_{-2}    & \cdots & a_{1-m}   & g_{0} (t) & g_{-1} (t) & \cdots & g_{-n} (t) \\
            c_{1}   & c_{0}     & c_{-1}    & \cdots & c_{2-m}   & b_{1}     & b_{0}      & \cdots & b_{1-n}    \\
            c_{2}   & c_{1}     & c_{0}     & \cdots & c_{3-m}   & b_{2}     & b_{1}      & \cdots & b_{2-n}    \\
            \vdots  & \vdots    & \vdots    & {}     & \vdots    & \vdots    & \vdots     & {}     & \vdots     \\
            c_{n}   & c_{n-1}   & c_{n-2}   & \cdots & c_{n+1-m} & b_{n}     & b_{n-1}    & \cdots & b_{0}      \\
            c_{n+1} & c_{n}     & c_{n-1}   & \cdots & c_{n+2-m} & b_{n+1}   & b_{n}      & \cdots & b_{1}      \\
            \vdots  & \vdots    & \vdots    & {}     & \vdots    & \vdots    & \vdots     & {}     & \vdots     \\
            c_{m+n} & c_{m+n-1} & c_{m+n-2} & \cdots & c_{n+1}   & b_{m+n}   & b_{m+n-1}  & \cdots & b_{m}      \\
        \end{bmatrix}.
    \end{align*}
    我们加行~\(1\) 的 \(t\) 倍于行~\(2\),
    再加行~\(2\) 的 \(t\) 倍于行~\(3\),
    ……
    再加行~\(m+n\) 的 \(t\) 倍于行~\(m+n+1\),
    得 \(m+n+1\)~级阵
    \begin{align*}
        S_1 =
        \begin{bmatrix}
            a_{0}   & a_{-1}    & a_{-2}    & \cdots & a_{1-m}   & g_{0} (t)   & g_{-1} (t)    & \cdots & g_{-n} (t)  \\
            a_{1}   & a_{0}     & a_{-1}    & \cdots & a_{2-m}   & g_{1} (t)   & g_{0} (t)     & \cdots & g_{1-n} (t) \\
            a_{2}   & a_{1}     & a_{0}     & \cdots & a_{3-m}   & g_{2} (t)   & g_{1} (t)     & \cdots & g_{2-n} (t) \\
            \vdots  & \vdots    & \vdots    & {}     & \vdots    & \vdots      & \vdots        & {}     & \vdots      \\
            a_{n}   & a_{n-1}   & a_{n-2}   & \cdots & a_{n+1-m} & g_{n} (t)   & g_{n-1} (t)   & \cdots & g_{0} (t)   \\
            a_{n+1} & a_{n}     & a_{n-1}   & \cdots & a_{n+2-m} & g_{n+1} (t) & g_{n} (t)     & \cdots & g_{1} (t)   \\
            \vdots  & \vdots    & \vdots    & {}     & \vdots    & \vdots      & \vdots        & {}     & \vdots      \\
            a_{m+n} & a_{m+n-1} & a_{m+n-2} & \cdots & a_{n+1}   & g_{m+n} (t) & g_{m+n-1} (t) & \cdots & g_{m} (t)   \\
        \end{bmatrix}.
    \end{align*}
    我们知道, \(\det {(S_1)} = \det {(S(h, g; n+1, m))}\).

    我们加列~\(m+2\) 的 \(-t\) 倍于列~\(m+1\),
    再加列~\(m+3\) 的 \(-t\) 倍于列~\(m+2\),
    ……
    再加列~\(m+n+1\) 的 \(-t\) 倍于列~\(m+n\),
    得 \(m+n+1\)~级阵
    \begin{align*}
        S_2 =
        \begin{bmatrix}
            a_{0}   & a_{-1}    & a_{-2}    & \cdots & a_{1-m}   & b_{0}   & b_{-1}    & \cdots & g_{-n} (t)  \\
            a_{1}   & a_{0}     & a_{-1}    & \cdots & a_{2-m}   & b_{1}   & b_{0}     & \cdots & g_{1-n} (t) \\
            a_{2}   & a_{1}     & a_{0}     & \cdots & a_{3-m}   & b_{2}   & b_{1}     & \cdots & g_{2-n} (t) \\
            \vdots  & \vdots    & \vdots    & {}     & \vdots    & \vdots  & \vdots    & {}     & \vdots      \\
            a_{n}   & a_{n-1}   & a_{n-2}   & \cdots & a_{n+1-m} & b_{n}   & b_{n-1}   & \cdots & g_{0} (t)   \\
            a_{n+1} & a_{n}     & a_{n-1}   & \cdots & a_{n+2-m} & b_{n+1} & b_{n}     & \cdots & g_{1} (t)   \\
            \vdots  & \vdots    & \vdots    & {}     & \vdots    & \vdots  & \vdots    & {}     & \vdots      \\
            a_{m+n} & a_{m+n-1} & a_{m+n-2} & \cdots & a_{n+1}   & b_{m+n} & b_{m+n-1} & \cdots & g_{m} (t)   \\
        \end{bmatrix}.
    \end{align*}
    我们知道, \(\det {(S_2)} = \det {(S_1)}\).
    则 \(\det {(S_2)} = \det {(S(h, g; n+1, m))}\).
    注意,
    \begin{align*}
        [S_2]_{i,j} =
        \begin{cases}
            a_{i-j},       & j \leq m;         \\
            b_{i-(j-m)},   & m < j \leq m + n; \\
            g_{i-1-n} (t), & j = m + n + 1.
        \end{cases}
    \end{align*}
    于是, 当 \(i = m + n + 1\) 时,
    若 \(j \leq m\),
    则 \([S_2]_{i,j} = a_{n+1+(m-j)} = 0\)
    (注意, \(n+1+(m-j) > n\)),
    且若 \(m < j \leq m + n\),
    则 \([S_2]_{i,j} = b_{m+1+(m+n-j)} = 0\)
    (注意, \(m+1+(m+n-j) > m\)).
    注意, \(S_2 ({m+n+1}|{m+n+1}) = S(f, g; n, m)\).
    注意, \(g_m (t) = g(t)\).
    于是, 按行~\(m+n+1\) 展开,
    \begin{align*}
        &
        \det {(S(h, g; n+1, m))}
        \\
        = {} &
        \det {(S_2)}
        \\
        = {} &
        (-1)^{(m+n+1) + (m+n+1)} g_m (t) \det {(S_2 ({m+n+1}|{m+n+1}))}
        \\
        = {} &
        \det {(S(f, g; n, m))}\, g(t).
        \qedhere
    \end{align*}
\end{proof}

\begin{theorem}
    设假定~\ref{standing-assumption-of-resultant} 成立.
    若 \(f(x) = a_0 (x - x_1) \dots (x - x_n)\), 则
    \begin{align*}
        \det {(S(f, g; n, m))} = a_0^m g(x_1) g(x_2) \dots g(x_n).
    \end{align*}
    类似地,
    若 \(g(x) = b_0 (x - y_1) \dots (x - y_m)\), 则
    \begin{align*}
        \det {(S(f, g; n, m))} = (-1)^{mn} b_0^n f(y_1) f(y_2) \dots f(y_m).
    \end{align*}
\end{theorem}

\begin{proof}
    我们证式~1;
    您用式~1 与定理~\ref{thm:resultant-swap} 证式~2.

    若 \(m = 0\),
    则式~1 的左边是 \(b_0^n\),
    且式~1 的右边是
    \(a_0^0 g(x_1) \dots g(x_n) = 1 \cdot b_0^n\)
    (注意, \(g(x) = b_0\)).

    以下, 我们设 \(m \geq 1\).

    记 \(F_0 (x) = a_0\),
    且 \(F_i (x) = F_{i-1} (x)\, (x - x_i)\).
    则
    \begin{align*}
        \det {(S(f, g; n, m))}
        = {} &
        \det {(S(F_n, g; n, m))}
        \\
        = {} &
        \det {(S(F_{n-1}, g; n-1, m))}\,
        g(x_n)
        \\
        = {} &
        \det {(S(F_{n-2}, g; n-2, m))}\,
        g(x_{n-1}) g(x_n)
        \\
        = {} &
        \dots
        \\
        = {} &
        \det {(S(F_0, g; 0, m))}\,
        g(x_1) g(x_2) \dots g(x_n)
        \\
        = {} &
        a_0^m g(x_1) g(x_2) \dots g(x_n).
        \qedhere
    \end{align*}
\end{proof}

\begin{theorem}
    设假定~\ref{standing-assumption-of-resultant} 成立.

    (1)
    设存在数 \(c\), 使 \(f(c) = 0 = g(c)\).
    则 \(\det {(S(f, g; n, m))} = 0\).

    (2)
    设 \(\det {(S(f, g; n, m))} = 0\).
    则 \(a_0 = 0 = b_0\),
    或当 \(a_0 \neq 0\) 或 \(b_0 \neq 0\) 时,
    存在数 \(c\), 使 \(f(c) = 0 = g(c)\).
\end{theorem}

\begin{proof}
    (1)
    我们其实已证它
    (定理~\ref{thm:necessary-condition-for-the-existence-of-a-common-root},
    定理~\ref{thm:resultant-Bézout}),
    但我们现在别地证它.

    若 \(a_0 = 0 = b_0\), 则 \(\det {(S(f, g; n, m))} = 0\).

    若 \(a_0 \neq 0\), 则我们可写
    \(f(x) = a_0 (x - x_1) (x - x_2) \dots (x - x_n)\),
    其中, \(x_1\), \(x_2\), \(\dots\), \(x_n\) 是数.
    因为 \(f(c) = 0\),
    \(a_0\), \(c - x_1\), \(c - x_2\), \(\dots\), \(c - x_n\)
    的一个是 \(0\).
    因为 \(a_0 \neq 0\),
    某 \(c - x_i\) 是 \(0\).
    则 \(c\) 等于此 \(x_i\).
    则
    \(\det {(S(f, g; n, m))}
    = a_0^m g(x_1) g(x_2) \dots g(x_n) = 0\),
    因为 \(g(x_i) = g(c) = 0\).

    若 \(b_0 \neq 0\), 您可类似地证它:
    注意,
    \begin{equation}
        \det {(S(f, g; n, m))}
        = (-1)^{m n} b_0^n f(y_1) f(y_2) \dots f(y_m),
        \label{eq:resultant-in-terms-of-roots-of-g}
    \end{equation}
    其中,
    \begin{equation}
        g(x) = b_0 (x - y_1) (x - y_2) \dots (x - y_m).
        \label{eq:roots-of-g}
    \end{equation}

    (2)
    若 \(a_0 = 0 = b_0\), 则 \(\det {(S(f, g; n, m))} = 0\).

    若 \(a_0 \neq 0\), 则我们可写
    \(f(x) = a_0 (x - x_1) (x - x_2) \dots (x - x_n)\),
    其中, \(x_1\), \(x_2\), \(\dots\), \(x_n\) 是数.
    注意,
    \(f(x_1)\), \(f(x_2)\), \(\dots\), \(f(x_n)\) 都是 \(0\).
    则
    \begin{align*}
        0
        = \det {(S(f, g; n, m))}
        = a_0^m g(x_1) g(x_2) \dots g(x_n).
    \end{align*}
    因为 \(a_0 \neq 0\),
    我们说,
    在 \(g(x_1)\), \(g(x_2)\), \(\dots\), \(g(x_n)\),
    \(0\) 存在.

    若 \(b_0 \neq 0\), 您可类似地证它:
    注意式~\eqref{eq:resultant-in-terms-of-roots-of-g},
    若式~\eqref{eq:roots-of-g}.
\end{proof}

\begin{theorem}
    设假定~\ref{standing-assumption-of-resultant} 成立.
    设 \(m \geq n\).

    设
    \begin{align*}
        q(x) = \sum_{k = 0}^{m-n} {c_k x^{m-n-k}}.
    \end{align*}
    为方便, 我们约定,
    若 \(k < 0\) 或 \(k > m - n\), 则 \(c_k = 0\).

    设
    \begin{align*}
        p(x) = g(x) + f(x) q(x) = \sum_{k = 0}^{m} {d_k x^{m-k}}.
    \end{align*}
    为方便, 我们约定,
    若 \(k < 0\) 或 \(k > m\), 则 \(d_k = 0\).

    则
    \begin{align*}
        & \det {(S(f, p; n, m))} = \det {(S(f, g; n, m))}, \\
        & \det {(S(p, f; m, n))} = \det {(S(g, f; m, n))}.
    \end{align*}
\end{theorem}

\begin{proof}
    我们证式~1;
    您用式~1 与定理~\ref{thm:resultant-swap} 证式~2.

    若 \(n = 0\),
    则 \(S(f, p; 0, m) = a_0 I_m = S(f, g; 0, m)\).
    % 则它们的行列式相等.

    以下, 我们设 \(n \geq 1\).

    由整式的加法与乘法,
    \begin{align*}
        d_k =
        \begin{dcases}
            b_k + \sum_{\ell = 0}^{k} {a_{k-\ell} c_{\ell}},
            & 0 \leq k \leq m; \\
            0,
            & \text{别的情形}.
        \end{dcases}
    \end{align*}
    注意, 若 \(0 \leq k \leq m\),
    \begin{align*}
        \sum_{\ell = 0}^{k} {a_{k-\ell} c_{\ell}}
        = \sum_{\ell = 0}^{m-n} {a_{k-\ell} c_{\ell}}.
    \end{align*}
    (若 \(k \leq m-n\), % tex-fmt: skip
    则当 \(\ell > k\) 时, \(a_{k-\ell} = 0\).
    若 \(k > m-n\),
    则当 \(\ell > m - n\) 时, \(c_{\ell} = 0\)). % tex-fmt: skip
    并且, 当 \(k < 0\) 或 \(k > m\) 时,
    \begin{align*}
        \sum_{\ell = 0}^{m-n} {a_{k-\ell} c_{\ell}} = 0.
    \end{align*}
    (若 \(k < 0\), % tex-fmt: skip
    则当 \(\ell \geq 0\) 时, \(k - \ell < 0\),
    故 \(a_{k-\ell} = 0\).
    若 \(k > m\),
    则当 \(\ell \leq m - n\) 时,
    \(k - \ell > m - \ell \geq m - (m - n) = n\),
    故 \(a_{k-\ell} = 0\).) % tex-fmt: skip
    于是, 统一地, 我们可写
    \begin{align*}
        d_k = b_k + \sum_{\ell = 0}^{m-n} {a_{k-\ell} c_{\ell}}.
    \end{align*}
    我们考虑 \(\det {(S(f, p; n, m))}\)
    与 \(\det {(S(f, g; n, m))}\)
    的关系.

    注意, \(m + n\) 级阵 \(S(f, g; n, m)\) 等于
    \begin{align*}
        \begin{bmatrix}
            a_{0}     & a_{-1}    & a_{-2}    & \cdots & a_{1-m}   & b_{0}     & b_{-1}    & \cdots & b_{1-n} \\
            a_{1}     & a_{0}     & a_{-1}    & \cdots & a_{2-m}   & b_{1}     & b_{0}     & \cdots & b_{2-n} \\
            a_{2}     & a_{1}     & a_{0}     & \cdots & a_{3-m}   & b_{2}     & b_{1}     & \cdots & b_{3-n} \\
            \vdots    & \vdots    & \vdots    & {}     & \vdots    & \vdots    & \vdots    & {}     & \vdots  \\
            a_{n-1}   & a_{n-2}   & a_{n-3}   & \cdots & a_{n-m}   & b_{n-1}   & b_{n-2}   & \cdots & b_{0}   \\
            a_{n}     & a_{n-1}   & a_{n-2}   & \cdots & a_{n+1-m} & b_{n}     & b_{n-1}   & \cdots & b_{1}   \\
            a_{n+1}   & a_{n}     & a_{n-1}   & \cdots & a_{n+2-m} & b_{n+1}   & b_{n}     & \cdots & b_{2}   \\
            \vdots    & \vdots    & \vdots    & {}     & \vdots    & \vdots    & \vdots    & {}     & \vdots  \\
            a_{m+n-1} & a_{m+n-2} & a_{m+n-3} & \cdots & a_{n}     & b_{m+n-1} & b_{m+n-2} & \cdots & b_{m}   \\
        \end{bmatrix},
    \end{align*}
    且 \(m + n\) 级阵 \(S(f, p; n, m)\) 等于
    \begin{align*}
        \begin{bmatrix}
            a_{0}     & a_{-1}    & a_{-2}    & \cdots & a_{1-m}   & d_{0}     & d_{-1}    & \cdots & d_{1-n} \\
            a_{1}     & a_{0}     & a_{-1}    & \cdots & a_{2-m}   & d_{1}     & d_{0}     & \cdots & d_{2-n} \\
            a_{2}     & a_{1}     & a_{0}     & \cdots & a_{3-m}   & d_{2}     & d_{1}     & \cdots & d_{3-n} \\
            \vdots    & \vdots    & \vdots    & {}     & \vdots    & \vdots    & \vdots    & {}     & \vdots  \\
            a_{n-1}   & a_{n-2}   & a_{n-3}   & \cdots & a_{n-m}   & d_{n-1}   & d_{n-2}   & \cdots & d_{0}   \\
            a_{n}     & a_{n-1}   & a_{n-2}   & \cdots & a_{n+1-m} & d_{n}     & d_{n-1}   & \cdots & d_{1}   \\
            a_{n+1}   & a_{n}     & a_{n-1}   & \cdots & a_{n+2-m} & d_{n+1}   & d_{n}     & \cdots & d_{2}   \\
            \vdots    & \vdots    & \vdots    & {}     & \vdots    & \vdots    & \vdots    & {}     & \vdots  \\
            a_{m+n-1} & a_{m+n-2} & a_{m+n-3} & \cdots & a_{n}     & d_{m+n-1} & d_{m+n-2} & \cdots & d_{m}   \\
        \end{bmatrix}.
    \end{align*}
    注意, \(j > m\) 时,
    \begin{align*}
        d_{i-(j-m)}
        = {} &
        b_{i-(j-m)} + \sum_{0 \leq \ell \leq m-n} {a_{i-(j-m)-\ell} c_{\ell}}
        \\
        = {} &
        b_{i-(j-m)} + \sum_{j-m \leq j-m+\ell \leq j-n} {a_{i-(j-m+\ell)} c_{(j-m+\ell)-(j-m)}}
        \\
        = {} &
        b_{i-(j-m)} + \sum_{j-m \leq r \leq j-m+(m-n)} {c_{r-(j-m)} a_{i-r}}.
    \end{align*}
    于是, 对 \(S(f, g; n, m)\), 我们%
    加列~\(1\) 的 \(c_0\)~倍于列~\(m+1\),
    再加列~\(2\) 的 \(c_1\)~倍于列~\(m+1\),
    ……
    再加列~\(m-n+1\) 的 \(c_{m-n}\)~倍于列~\(m+1\),
    且我们加列~\(2\) 的 \(c_0\)~倍于列~\(m+2\),
    再加列~\(3\) 的 \(c_1\)~倍于列~\(m+2\),
    ……
    再加列~\(m-n+2\) 的 \(c_{m-n}\)~倍于列~\(m+2\),
    ……
    且我们加列~\(n\) 的 \(c_0\)~倍于列~\(m+n\),
    再加列~\(n+1\) 的 \(c_1\)~倍于列~\(m+n\),
    ……
    再加列~\(m\) 的 \(c_{m-n}\)~倍于列~\(m+n\),
    则我们得 \(S(f, p; n, m)\).
    则
    \begin{equation*}
        \det {(S(f, p; n, m))} = \det {(S(f, g; n, m))}.
        \qedhere
    \end{equation*}
\end{proof}

\begin{example}
    设假定~\ref{standing-assumption-of-resultant} 成立.
    设
    \begin{align*}
        G(x) = \sum_{k = 0}^{m+1} {c_k x^{m+1-k}},
    \end{align*}
    其中, \(c_k = b_{k-1}\);
    特别地, \(c_0 = 0\).
    注意, 其实, \(G(x) = g(x)\),
    但
    \begin{align*}
        \sum_{k = 0}^{m} {b_k x^{m-k}}
        \quad \text{与} \quad
        \sum_{k = 0}^{m+1} {c_k x^{m+1-k}}
    \end{align*}
    的形式是不同的.
    \(S(f, g; n, m)\) 与 \(S(f, G; n, m+1)\)
    分别是 \(m+n\)~级与 \(m+n+1\)~级阵.
    我们考虑 \(\det {(S(f, g; n, m))}\)
    与 \(\det {(S(f, G; n, m+1))}\)
    的关系.

    首先, 既然 \(c_0 = 0\),
    则 \(k \leq 0\) 时, \(c_k \leq 0\).
    则 \(j > 1\) 时, \([S(f, G; n, m+1)]_{1,j} = 0\).
    (若 \(2 \leq j \leq m+1\), % tex-fmt: skip
    则 \([S(f, G; n, m+1)]_{1,j} = a_{1-j} = 0\);
    若 \(j \geq m+2\), 则
    \([S(f, G; n, m+1)]_{1,j} % tex-fmt: skip
    = c_{1-(j-(m+1))} = c_{m+2-j} = 0\).) % tex-fmt: skip

    另一方面,
    \((S(f, G; n, m+1))(1|1) = S(f, g; n, m)\).
    注意,
    \begin{align*}
        [S(f, G; n, m+1)]_{i+1,j+1}
        = {} &
        \begin{cases}
            a_{(i+1)-(j+1)},       & j+1 \leq m+1; \\
            c_{(i+1)-(j+1-(m+1))}, & j+1 > m+1
        \end{cases}
        \\
        = {} &
        \begin{cases}
            a_{i-j},       & j \leq m; \\
            c_{i-(j-m)+1}, & j > m
        \end{cases}
        \\
        = {} &
        \begin{cases}
            a_{i-j},     & j \leq m; \\
            b_{i-(j-m)}, & j > m
        \end{cases}
        \\
        = {} &
        [S(f, g; n, m)]_{i,j}.
    \end{align*}
    所以, 由按行~\(1\) 展开,
    \begin{align*}
        &
        \det {(S(f, G; n, m+1))}
        \\
        = {} &
        (-1)^{1+1}\, [S(f, G; n, m+1)]_{1,1}
        \det {((S(f, G; n, m+1))(1|1))}
        \\
        = {} &
        a_0 \det {(S(f, g; n, m))}.
    \end{align*}

    注意, \(g = G\).
    为方便, 我们直接地写
    \begin{align*}
        \det {(S(f, g; n, m+1))} = a_0 \det {(S(f, g; n, m))}.
    \end{align*}
    由此, 若 \(d\) 是非负整数,
    \begin{align*}
        \det {(S(f, g; n, m+d))}
        = {} & a_0 \det {(S(f, g; n, m+d-1))}   \\
        = {} & a_0^2 \det {(S(f, g; n, m+d-2))} \\
        = {} & \dots                            \\
        = {} & a_0^d \det {(S(f, g; n, m))},
    \end{align*}
    其中, \(S(f, g; n, m+d)\) 的 \(g\) 被认为是
    \begin{align*}
        g(x) = \sum_{k = 0}^{m+d} {w_k x^{m+d-k}},
    \end{align*}
    其中, \(w_k = b_{k-d}\).

    最后, 类似地,
    \begin{align*}
        \det {(S(f, g; n+d, m))}
        = {} &
        (-1)^{m (n+d)} \det {(S(g, f; m, n+d))}
        \\
        = {} &
        (-1)^{m n} (-1)^{m d} b_0^d \det {(S(g, f; m, n))}
        \\
        = {} &
        (-1)^{m d} b_0^d \det {(S(f, g; n, m))}.
        \qefhere
    \end{align*}
\end{example}

\begin{example}
    设假定~\ref{standing-assumption-of-resultant} 成立.
    设
    \begin{align*}
        F(x) = \sum_{k = 0}^{n} {c_k x^{n-k}},
        \quad
        G(x) = \sum_{k = 0}^{m} {d_k x^{m-k}},
    \end{align*}
    其中, (对任何整数 \(\ell\),)
    \(c_\ell = a_{n-\ell}\),
    且 \(d_\ell = b_{m-\ell}\).
    于是,
    \begin{align*}
        F(x) & = a_n x^n + a_{n-1} x^{n-1} + \dots + a_0, \\
        G(x) & = b_m x^m + b_{m-1} x^{m-1} + \dots + b_0.
    \end{align*}
    我们考虑
    \(\det {(S(F, G; n, m))}\)
    与
    \(\det {(S(f, g; n, m))}\)
    的关系.

    注意, 若 \(j \leq m\), 则
    \begin{align*}
        [S(F, G; n, m)]_{i,j}
        = {} & c_{i-j}                          \\
        = {} & a_{n-(i-j)}                      \\
        = {} & a_{n-i+j}                        \\
        = {} & a_{m+n+1-i+j-m-1}                \\
        = {} & a_{(m+n+1-i)-(m+1-j)}            \\
        = {} & [S(f, g; n, m)]_{m+n+1-i,m+1-j};
    \end{align*}
    若 \(j > m\), 则
    \begin{align*}
        [S(F, G; n, m)]_{i,j}
        = {} & d_{i-(j-m)}                            \\
        = {} & b_{m-(i-(j-m))}                        \\
        = {} & b_{m-i+(j-m)}                          \\
        = {} & b_{m+n+1-i+(j-m)-n-1}                  \\
        = {} & b_{(m+n+1-i)-(n+1-(j-m))}              \\
        = {} & b_{(m+n+1-i)-(m+n+1-(j-m)-m)}          \\
        = {} & [S(f, g; n, m)]_{m+n+1-i,m+n+1-(j-m)}.
    \end{align*}
    作 \(m+n\)~级阵 \(T\), 使
    \([T]_{i,j} = [S(f, g; n, m)]_{m+n+1-i,j}\).
    则 \(T\) 的行~\(i\) 是 \(S(f, g; n, m)\) 的行~\(m+n+1-i\).
    注意,
    \begin{align*}
        [S(F, G; n, m)]_{i,j} =
        \begin{cases}
            [T]_{i,m+1-j},       & j \leq m; \\
            [T]_{i,m+n+1-(j-m)}, & j > m.
        \end{cases}
    \end{align*}
    则 \(T\) 的列~\(j\) 是 \(S(F, G; n, m)\) 的列~\(m+1-j\)
    (若 \(j \leq m\))
    或列~\(m+n+1-(j-m)\)
    (若 \(j > m\)).
    我们分别考虑
    \(\det {(T)}\), \(\det {(S(f, g; n, m))}\) 的关系,
    与
    \(\det {(T)}\), \(\det {(S(F, G; n, m))}\)
    的关系.

    对 \(S(f, g; n, m)\),
    我们交换行~\(m+n\), \(m+n-1\),
    再交换行~\(m+n-1\), \(m+n-2\),
    ……
    再交换行~\(2\), \(1\),
    得 \(m+n\)~级阵 \(T_1\).
    则 \(T_1\) 的行~\(1\) 是 \(S(f, g; n, m)\) 的行~\(m+n\),
    且 \(T_1\) 的行~\(\ell\) 是 \(S(f, g; n, m)\) 的行~\(\ell-1\)
    (若 \(\ell > 1\)).
    我们作了 \(m+n-1\) 次相邻的行的交换.
    则 \(\det {(T_1)} = (-1)^{m+n-1} \det {(S(f, g; n, m))}\).
    注意, \(T_1\) 的行~\(1\) 是 \(T\) 的行~\(1\).

    我们可类似地使 \(T_1\)~的行~\(m+n\)
    (其是 \(S(f, g; n, m)\) 的行~\(m+n-1\))
    到行~\(2\) 的位置,
    使 \(T_1\) 的行~\(1\) 仍在行~\(1\) 的位置,
    且使 \(T_1\) 的行~\(\ell\) 在行~\(\ell+1\) 的位置
    (若 \(2 < \ell < m+n\)),
    得 \(m+n\)~级阵 \(T_2\).
    我们作了 \(m+n-2\) 次相邻的行的交换.
    则
    \begin{align*}
        \det {(T_2)}
        = {} & (-1)^{m+n-2} \det {(T_1)}                      \\
        = {} & (-1)^{(m+n-1)+(m+n-2)} \det {(S(f, g; n, m))}.
    \end{align*}

    ……

    最后, 我们共作
    \begin{align*}
        (m+n-1) + (m+n-2) + \dots + 1
        = {} &
        \frac{(m+n)(m+n-1)}{2}
        \\
        = {} &
        \frac{m (m-1+n) + n (n-1+m)}{2}
        \\
        = {} &
        \frac{m(m-1) + n(n-1) + 2mn}{2}
        \\
        = {} &
        \frac{m(m-1)}{2} + \frac{n(n-1)}{2} + mn
    \end{align*}
    次相邻的行的交换,
    使 \(S(f, g; n, m)\) 的行~\(\ell\)
    到行~\(m+n-1-\ell\) 的位置,
    \(\ell = 1\), \(2\), \(\dots\), \(m+n\),
    得 \(m+n\)~级阵 \(T\).
    则
    \begin{align*}
        \det {(T)}
        = (-1)^{m(m-1)/2} (-1)^{n(n-1)/2} (-1)^{mn}
        \det {(S(f, g; n, m))}.
    \end{align*}

    然后, 我们可类似地,
    作 \(m + (m-1) + \dots + 1 = m(m-1)/2\)
    次相邻的列的交换,
    且再作 \(n + (n-1) + \dots + 1 = n(n-1)/2\)
    次相邻的列的交换,
    使 \(T\) 的列~\(1\), \(2\), \(\dots\), \(m\)
    分别地到列~\(m\), \(m-1\), \(\dots\), \(1\) 的位置,
    且使 \(T\) 的列~\(m+1\), \(m+2\), \(\dots\), \(m+n\)
    分别地到列~\(m+n\), \(m+n-1\), \(\dots\), \(m+1\) 的位置,
    得 \(m+n\)~级阵 \(S(F, G; n, m)\).
    则
    \begin{align*}
        &
        \det {(S(F, G; n, m))}
        \\
        = {} &
        (-1)^{n(n-1)/2} (-1)^{m(m-1)/2} \det {(T)}
        \\
        = {} &
        (-1)^{n(n-1)/2} (-1)^{m(m-1)/2}
        (-1)^{m(m-1)/2} (-1)^{n(n-1)/2} (-1)^{mn}
        \det {(S(f, g; n, m))}
        \\
        = {} &
        (-1)^{mn} \det {(S(f, g; n, m))}.
    \end{align*}

    我们以 \(a_\ell\), \(b_\ell\)
    直接地表示 \(S(F, G; n, m)\) 的元:
    \begin{align*}
        [S(F, G; n, m)]_{i,j}
        = {} &
        \begin{cases}
            a_{n-(i-j)},     & j \leq m; \\
            b_{m-(i-(j-m))}, & j > m;
        \end{cases}
        \\
        = {} &
        \begin{cases}
            a_{j-i+n}, & j \leq m; \\
            b_{j-i},   & j > m.
        \end{cases}
    \end{align*}
    对适合假定~\ref{standing-assumption-of-resultant}
    的 \(f\), \(g\),
    定义 \(m + n\)~级阵 \(S'(f, g; n, m)\),
    其中,
    \begin{align*}
        [S'(f, g; n, m)]_{i,j} =
        \begin{cases}
            a_{j-i+n}, & j \leq m; \\
            b_{j-i},   & j > m.
        \end{cases}
    \end{align*}
    比如, 若 \(n = 2\), 且 \(m = 5\),
    则 \(S'(f, g; 2, 5)\) 等于
    \begin{align*}
        \begin{bmatrix}
            a_2    & a_3    & a_4    & a_5    & a_6 & b_5    & b_6 \\
            a_1    & a_2    & a_5    & a_4    & a_5 & b_4    & b_5 \\
            a_0    & a_1    & a_2    & a_3    & a_4 & b_3    & b_4 \\
            a_{-1} & a_0    & a_1    & a_2    & a_3 & b_2    & b_3 \\
            a_{-2} & a_{-1} & a_0    & a_1    & a_2 & b_1    & b_2 \\
            a_{-3} & a_{-2} & a_{-1} & a_0    & a_1 & b_0    & b_1 \\
            a_{-4} & a_{-3} & a_{-2} & a_{-1} & a_0 & b_{-1} & b_0 \\
        \end{bmatrix}
        =
        \begin{bmatrix}
            a_2 & 0   & 0   & 0   & 0   & b_5 & 0   \\
            a_1 & a_2 & 0   & 0   & 0   & b_4 & b_5 \\
            a_0 & a_1 & a_2 & 0   & 0   & b_3 & b_4 \\
            0   & a_0 & a_1 & a_2 & 0   & b_2 & b_3 \\
            0   & 0   & a_0 & a_1 & a_2 & b_1 & b_2 \\
            0   & 0   & 0   & a_0 & a_1 & b_0 & b_1 \\
            0   & 0   & 0   & 0   & a_0 & 0   & b_0 \\
        \end{bmatrix}.
    \end{align*}

    在一些文献, 一些作者的
    ``\(f\), \(g\) 的 Sylvester 阵'' 是这儿的
    \(S'(f, g; n, m)\),
    且 ``\(f\), \(g\) 的结式'' 是
    \(\det {(S'(f, g; n, m))}\).
    % 形式地, \(S(f, g; n, m)\) 是 ``降幂的''
    % (\(a_0\), \(a_1\), \(\dots\), \(a_n\)
    % 分别是 \(f(x)\) 的
    % \(x^n\), \(x^{n-1}\), \(\dots\), \(1\) 项的系数,
    % 且 \(b_0\), \(b_1\), \(\dots\), \(b_m\)
    % 分别是 \(g(x)\) 的
    % \(x^m\), \(x^{m-1}\), \(\dots\), \(1\) 项的系数),
    % 而 \(S'(f, g; n, m)\) 是 ``升幂的''
    % (\(a_n\), \(a_{n-1}\), \(\dots\), \(a_0\)
    % 分别是 \(f(x)\) 的
    % \(1\), \(x\), \(\dots\), \(x^n\) 项的系数,
    % 且 \(b_m\), \(b_{m-1}\), \(\dots\), \(b_0\)
    % 分别是 \(g(x)\) 的
    % \(1\), \(x\), \(\dots\), \(x^m\) 项的系数).
    % 由前面的讨论, 它们的行列式
    % (即 \(f\), \(g\) 的 \(2\)~类结式)
    由前面的讨论, 它们的行列式%
    至多差正负号:
    \(\det {(S'(f, g; n, m))}
    = (-1)^{m n} \det {(S(f, g; n, m))}\).
\end{example}

\section{derivative}

本节, 我们讨论整式的 derivative.
这会是有用的工具.

\begin{definition}
    设 \(n\) 是非负整数.
    设
    \begin{align*}
        f(x) = a_0 x^n + a_1 x^{n-1} + \dots + a_{n-1} x + a_n
        = \sum_{i = 0}^{n} {a_i x^{n-i}}.
    \end{align*}
    是整式.
    定义它的 \emph{derivative} 为整式
    \begin{align*}
        \mathrm{D} f(x)
        = n a_0 x^{n-1} + (n-1) a_1 x^{n-2} + \dots + 1 a_{n-1}
        = \sum_{i = 0}^{n-1} {(n-i) a_i x^{n-1-i}}.
    \end{align*}
\end{definition}

\begin{example}
    设 \(f(x) = a x^2 + b x + c\).
    则 \(\mathrm{D} f(x) = 2 a x + b\).
\end{example}

\begin{example}
    设 \(f(x) = c x^i\) (\(i \geq 1\)).
    我们可视 \(f(x)\) 为 \(c x^i + 0 x^{i-1} + \dots + 0\).
    则 \(\mathrm{D} f(x)
        = i c x^{i-1} + (i-1) 0 x^{i-2} + \dots + 0
    = i c x^{i-1}\).

    值得提, \(\mathrm{D} (c x^0) = 0\).
\end{example}

\begin{theorem}
    设 \(r\) 是正整数.
    设 \(f_1 (x)\), \(f_2 (x)\), \(\dots\), \(f_r (x)\) 是整式.
    设 \(c_1\), \(c_2\), \(\dots\), \(c_m\) 是数.
    则
    \begin{align*}
        \mathrm{D} \sum_{k = 1}^{r} {c_k f_k (x)}
        = \sum_{k = 1}^{r} {c_k \mathrm{D} f_k (x)}.
    \end{align*}
\end{theorem}

\begin{proof}
    我们设, 非负整数 \(n\) 适合,
    对每个不超过 \(r\) 的正整数 \(k\),
    \begin{align*}
        f_k (x)
        = a_{k,0} x^n + a_{k,1} x^{n-1}
        + \dots + a_{k,n-1} x + a_{k,n}.
    \end{align*}
    (注意, \(a_{k,0}\) 可以为 \(0\).)
    则
    \begin{align*}
        \mathrm{D} \sum_{k = 1}^{r} {c_k f_k (x)}
        = {} &
        \mathrm{D}
        \sum_{k = 1}^{r}
        c_k
        \sum_{i = 0}^{n}
        a_{k,i} x^{n-i}
        \\
        = {} &
        \mathrm{D}
        \sum_{k = 1}^{r}
        \sum_{i = 0}^{n}
        c_k a_{k,i} x^{n-i}
        \\
        = {} &
        \mathrm{D}
        \sum_{i = 0}^{n}
        \sum_{k = 1}^{r}
        c_k a_{k,i} x^{n-i}
        \\
        = {} &
        \mathrm{D}
        \sum_{i = 0}^{n}
        \left(\sum_{k = 1}^{r} c_k a_{k,i}\right)
        x^{n-i}
        \\
        = {} &
        \sum_{i = 0}^{n-1}
        (n-i) \left(\sum_{k = 1}^{r} c_k a_{k,i}\right)
        x^{n-1-i}
        \\
        = {} &
        \sum_{i = 0}^{n-1}
        \sum_{k = 1}^{r}
        (n-i) c_k a_{k,i} x^{n-1-i}
        \\
        = {} &
        \sum_{k = 1}^{r}
        \sum_{i = 0}^{n-1}
        (n-i) c_k a_{k,i} x^{n-1-i}
        \\
        = {} &
        \sum_{k = 1}^{r}
        c_k
        \sum_{i = 0}^{n-1}
        (n-i) a_{k,i} x^{n-1-i}
        \\
        = {} &
        \sum_{k = 1}^{r}
        c_k
        \mathrm{D} f_k (x).
        \qedhere
    \end{align*}
\end{proof}

\begin{theorem}
    设 \(f(x)\), \(g(x)\) 是整式.
    则
    \begin{align*}
        \mathrm{D} (f(x) g(x))
        = (\mathrm{D} f(x)) g(x) + f(x) (\mathrm{D} g(x)).
    \end{align*}
\end{theorem}

\begin{proof}
    我们设
    \begin{align*}
        f(x) & = a_0 x^n + a_1 x^{n-1} + \dots + a_{n-1} x + a_n, \\
        g(x) & = b_0 x^m + b_1 x^{m-1} + \dots + b_{m-1} x + b_m.
    \end{align*}

    我们先证明, 对任何非负整数 \(i\), \(j\),
    与任何数 \(p\), \(q\),
    \begin{align*}
        \mathrm{D} (p x^i \cdot q x^j)
        = \mathrm{D} (p x^i) \cdot q x^j
        + p x^i \cdot \mathrm{D} (q x^j).
    \end{align*}
    若 \(i = 0 = j\), 则
    \begin{align*}
        \mathrm{D} (p x^i \cdot q x^j)
        = {} &
        \mathrm{D} (p q)
        \\
        = {} &
        0
        \\
        = {} &
        0 \cdot q x^j + p x^i \cdot 0
        \\
        = {} &
        \mathrm{D} (p x^i) \cdot q x^j
        + p x^i \cdot \mathrm{D} (q x^j).
    \end{align*}
    若 \(i = 0\), 且 \(j \geq 1\), 则
    \begin{align*}
        \mathrm{D} (p x^i \cdot q x^j)
        = {} &
        \mathrm{D} (p q x^j)
        \\
        = {} &
        j p q x^{j-1}
        \\
        = {} &
        0 \cdot q x^j + p x^i \cdot j q x^{j-1}
        \\
        = {} &
        \mathrm{D} (p x^i) \cdot q x^j
        + p x^i \cdot \mathrm{D} (q x^j).
    \end{align*}
    若 \(i \geq 1\), 且 \(j = 0\), 则
    \begin{align*}
        \mathrm{D} (p x^i \cdot q x^j)
        = {} &
        \mathrm{D} (p q x^i)
        \\
        = {} &
        i p q x^{i-1}
        \\
        = {} &
        i p x^{i-1} \cdot q x^j + p x^i \cdot 0
        \\
        = {} &
        \mathrm{D} (p x^i) \cdot q x^j
        + p x^i \cdot \mathrm{D} (q x^j).
    \end{align*}
    若 \(i \geq 1\), 且 \(j \geq 1\), 则
    \begin{align*}
        \mathrm{D} (p x^i \cdot q x^j)
        = {} &
        \mathrm{D} (p q x^{i+j})
        \\
        = {} &
        (i + j) p q x^{i+j-1}
        \\
        = {} &
        i p q x^{i+j-1} + j p q x^{i+j-1}
        \\
        = {} &
        i p x^{i-1} \cdot q x^j + p x^i \cdot j q x^{j-1}
        \\
        = {} &
        \mathrm{D} (p x^i) \cdot q x^j
        + p x^i \cdot \mathrm{D} (q x^j).
    \end{align*}

    由此, 且由上个定理,
    \begin{align*}
        \mathrm{D} (f(x) g(x))
        = {} &
        \mathrm{D} \left(
            \left(\sum_{i=0}^{n} {a_i x^{n-i}}\right)
            \left(\sum_{j=0}^{m} {b_j x^{m-j}}\right)
        \right)
        \\
        = {} &
        \mathrm{D} \left(
            \sum_{i=0}^{n} \sum_{j=0}^{m}
            {a_i x^{n-i} \cdot b_j x^{m-j}}
        \right)
        \\
        = {} &
        \sum_{i=0}^{n} \sum_{j=0}^{m}
        {\mathrm{D} (a_i x^{n-i} \cdot b_j x^{m-j})}
        \\
        = {} &
        \sum_{i=0}^{n} \sum_{j=0}^{m}
        {(\mathrm{D} (a_i x^{n-i}) \cdot b_j x^{m-j}
        + a_i x^{n-i} \cdot \mathrm{D} (b_j x^{m-j}))}
        \\
        = {} &
        \sum_{i=0}^{n}
        \left(
            \mathrm{D} (a_i x^{n-i}) \sum_{j=0}^{m} b_j x^{m-j}
            + a_i x^{n-i} \sum_{j=0}^{m} \mathrm{D} (b_j x^{m-j})
        \right)
        \\
        = {} &
        \sum_{i=0}^{n}
        \left(
            \mathrm{D} (a_i x^{n-i}) \sum_{j=0}^{m} b_j x^{m-j}
            + a_i x^{n-i} \mathrm{D} \sum_{j=0}^{m} b_j x^{m-j}
        \right)
        \\
        = {} &
        \sum_{i=0}^{n}
        ( \mathrm{D} (a_i x^{n-i}) g(x)
        + a_i x^{n-i} \mathrm{D} g(x))
        \\
        = {} &
        \left(
            \sum_{i=0}^{n} \mathrm{D} (a_i x^{n-i})
        \right) g(x)
        + \left(
            \sum_{i=0}^{n} a_i x^{n-i}
        \right) \mathrm{D} g(x)
        \\
        = {} &
        \left(\mathrm{D} \sum_{i=0}^{n} a_i x^{n-i}\right)
        g(x)
        + f(x) (\mathrm{D} g(x))
        \\
        = {} &
        (\mathrm{D} f(x)) g(x) + f(x) (\mathrm{D} g(x)).
        \qedhere
    \end{align*}
\end{proof}

\begin{theorem}
    设 \(r\) 是正整数.
    设 \(f_1 (x)\), \(f_2 (x)\), \(\dots\), \(f_r (x)\) 是整式.
    则
    \begin{align*}
        \mathrm{D} (f_1 (x) f_2 (x) \dots f_r (x))
        = {} &
        (\mathrm{D} f_1 (x)) f_2 (x) \dots f_r (x)
        + f_1 (x) (\mathrm{D} f_2 (x)) \dots f_r (x)
        \\
        &
        \quad
        + \dots
        + f_1 (x) f_2 (x) \dots (\mathrm{D} f_r (x)).
    \end{align*}
    特别地, 若 \(f_1 (x)\), \(f_2 (x)\), \(\dots\), \(f_r (x)\)
    都是 \(f(x)\),
    则
    \begin{align*}
        \mathrm{D} ((f(x))^r)
        = r (f(x))^{r-1} \mathrm{D} f(x).
    \end{align*}
\end{theorem}

\begin{proof}
    作命题 \(P(r)\):
    对任何整式
    \(f_1 (x)\), \(f_2 (x)\), \(\dots\), \(f_r (x)\),
    \begin{align*}
        \mathrm{D} (f_1 (x) f_2 (x) \dots f_r (x))
        = {} &
        (\mathrm{D} f_1 (x)) f_2 (x) \dots f_r (x)
        + f_1 (x) (\mathrm{D} f_2 (x)) \dots f_r (x)
        \\
        &
        \quad
        + \dots
        + f_1 (x) f_2 (x) \dots (\mathrm{D} f_r (x)).
    \end{align*}
    我们用数学归纳法证明,
    对任何正整数 \(r\), \(P(r)\) 是对的.

    \(P(1)\) 是对的, 因为左边与右边都是 \(D f_1 (x)\).

    设 \(P(r)\) 是对的.
    我们由此证明, \(P(r+1)\) 是对的.
    注意, 由假定,
    \begin{align*}
        &
        \mathrm{D} (f_1 (x) f_2 (x) \dots f_{r+1} (x))
        \\
        = {} &
        \mathrm{D} ((f_1 (x) f_2 (x) \dots f_r (x)) f_{r+1} (x))
        \\
        = {} &
        (\mathrm{D} (f_1 (x) f_2 (x) \dots f_r (x))) f_{r+1} (x)
        + (f_1 (x) f_2 (x) \dots f_r (x)) (\mathrm{D} f_{r+1} (x))
        \\
        = {} &
        ((\mathrm{D} f_1 (x)) f_2 (x) \dots f_r (x)
            + f_1 (x) (\mathrm{D} f_2 (x)) \dots f_r (x)
            \\
            &
            \quad
            + \dots
        + f_1 (x) f_2 (x) \dots (\mathrm{D} f_r (x))) f_{r+1} (x)
        + (f_1 (x) f_2 (x) \dots f_r (x)) (\mathrm{D} f_{r+1} (x))
        \\
        = {} &
        (\mathrm{D} f_1 (x)) f_2 (x) \dots f_r (x) f_{r+1} (x)
        + f_1 (x) (\mathrm{D} f_2 (x)) \dots f_r (x) f_{r+1} (x)
        \\
        &
        \quad
        + \dots
        + f_1 (x) f_2 (x) \dots (\mathrm{D} f_r (x)) f_{r+1} (x)
        + f_1 (x) f_2 (x) \dots f_r (x) (\mathrm{D} f_{r+1} (x)).
    \end{align*}
    所以, \(P(r+1)\) 是对的.
    由数学归纳法, 待证的命题成立.
\end{proof}

\begin{theorem}
    设 \(f(x) = (x - c) q(x)\).
    则 \(\mathrm{D} f(c) = q(c)\).

    若 \(f(x) = a_0 (x - x_1) (x - x_2) \dots (x - x_n)\),
    则对不超过 \(n\) 的正整数 \(i\),
    \begin{align*}
        \mathrm{D} f(x_i) = a_0 \prod_{\substack{1 \leq k \leq n \\ k \neq i}}
        {(x_i - x_k)}.
    \end{align*}
\end{theorem}

\begin{proof}
    注意,
    \begin{align*}
        \mathrm{D} f(x)
        = \mathrm{D} (x - x_0) q(x) + (x - x_0) \mathrm{D} q(x)
        = q(x) + (x - x_0) \mathrm{D} q(x).
    \end{align*}
    代 \(x\) 以 \(c\), 得
    \begin{align*}
        \mathrm{D} f(c)
        = q(c) + (c - c) \mathrm{D} q(c)
        = q(c).
    \end{align*}

    若 \(f(x) = a_0 (x - x_1) (x - x_2) \dots (x - x_n)\),
    则取
    \begin{align*}
        q_i (x)
        = a_0 \prod_{\substack{1 \leq k \leq n \\ k \neq i}} {(x - x_k)},
    \end{align*}
    得 \(f(x) = (x - x_i) q_i (x)\).
    则
    \begin{equation*}
        \mathrm{D} f(x_i) = q_i (x_i)
        = a_0 \prod_{\substack{1 \leq k \leq n \\ k \neq i}} {(x_i - x_k)}.
        \qedhere
    \end{equation*}
\end{proof}

最后, 我们以整式的几个性质结束本节.

\begin{theorem}
    设 \(n\) 是非负整数.
    设 \(f(x) = a_0 x^n + a_1 x^{n-1} + \dots + a_n\) 是整式.
    设 \(c\) 是数.
    则存在整式 \(q(x)\), 使
    \begin{align*}
        f(x) = (x - c) q(x) + f(c).
    \end{align*}
    进一步地, 若 \(n \geq 1\), 且 \(a_0 \neq 0\),
    则 \(q(x)\) 的次是 \(n-1\).
\end{theorem}

\begin{proof}
    对任何正整数 \(m\), 记
    \begin{align*}
        s_m (x; y)
        = x^{m-1} + x^{m-2} y + \dots + x y^{m-2} + y^{m-1}
        = \sum_{k = 0}^{m-1} {x^{m-1-k} y^k}.
    \end{align*}
    则 \(s_m (x; y)\) 是关于 \(x\) 的 \(m-1\)~次整式, 且
    \begin{align*}
        (x - y) s_m (x; y)
        = {} &
        (x - y) \sum_{k = 0}^{m-1} {x^{m-1-k} y^k}
        \\
        = {} &
        x \sum_{k = 0}^{m-1} {x^{m-1-k} y^k}
        - y \sum_{k = 0}^{m-1} {x^{m-1-k} y^k}
        \\
        = {} &
        \sum_{k = 0}^{m-1} {x^{m-k} y^k}
        - \sum_{k = 0}^{m-1} {x^{m-1-k} y^{k+1}}
        \\
        = {} &
        \sum_{\ell = 0}^{m-1} {x^{m-\ell} y^\ell}
        - \sum_{\ell = 1}^{m} {x^{m-\ell} y^\ell}
        \\
        = {} &
        x^{m-0} y^0 + \sum_{\ell = 1}^{m-1} {x^{m-\ell} y^\ell}
        - \left(
            \sum_{\ell = 1}^{m-1} {x^{m-\ell} y^\ell} + x^{m-m} y^m
        \right)
        \\
        = {} &
        x^m - y^m.
    \end{align*}
    则
    \begin{align*}
        f(x) - f(c)
        = {} &
        \sum_{i = 0}^{n} {a_i x^{n-i}}
        - \sum_{i = 0}^{n} {a_i c^{n-i}}
        \\
        = {} &
        \sum_{i = 0}^{n} {(a_i x^{n-i} - a_i c^{n-i})}
        \\
        = {} &
        \sum_{i = 0}^{n} {a_i (x^{n-i} - c^{n-i})}
        \\
        = {} &
        \sum_{i = 0}^{n-1} {a_i (x^{n-i} - c^{n-i})}
        + a_n (x^0 - c^0)
        \\
        = {} &
        \sum_{i = 0}^{n-1} {a_i (x - c) s_{n-i} (x; c)}
        + a_n\, 0
        \\
        = {} &
        (x - c)
        \sum_{i = 0}^{n-1} {a_i s_{n-i} (x; c)}.
    \end{align*}
    记
    \begin{align*}
        q(x) = \sum_{i = 0}^{n-1} {a_i s_{n-i} (x; c)}.
    \end{align*}
    则 \(f(x) = (x - c) q(x) + f(c)\).

    进一步地, 若 \(n \geq 1\), 则
    \begin{align*}
        q(x)
        = {} &
        a_0 s_n (x; c)
        + \underbrace{\sum_{i = 1}^{n-1} {a_i s_{n-i} (x; c)}}_{\text{次低于 \(n-1\)}}
        \\
        = {} &
        a_0 (x^{n-1} + c x^{n-2} + \dots + c^{n-1})
        + (\text{次低于 \(n-1\) 的整式})
        \\
        = {} &
        a_0 x^{n-1} + (\text{次低于 \(n-1\) 的整式}).
    \end{align*}
    因为 \(a_0 \neq 0\), \(q(x)\) 的次是 \(n-1\).
\end{proof}

\begin{theorem}
    设 \(n\) 是正整数.
    设 \(f(x) = a_0 x^n + a_1 x^{n-1} + \dots + a_n\)
    (\(a_0 \neq 0\)) 是 \(n\)~次整式.

    (1)
    若存在数 \(d\), 与 \(n-1\)~次整式 \(p(x)\),
    使 \(f(x) = (x - d) p(x)\),
    则 \(f(d) = 0\).

    (2)
    若数 \(c\) 适合, \(f(c) = 0\),
    则存在 \(n-1\)~次整式 \(q(x)\),
    使 \(f(x) = (x - c) q(x)\).
\end{theorem}

\begin{proof}
    (1)
    注意, \(f(d) = (d - d) p(d) = 0\, p(d) = 0\).

    (2)
    由上个定理, 存在 \(n-1\) 次整式 \(q(x)\), 使
    \begin{equation*}
        f(x) = (x - c) q(x) + f(c)
        = (x - c) q(x) + 0
        = (x - c) q(x).
        \qedhere
    \end{equation*}
\end{proof}

\begin{theorem}
    设 \(n\) 是正整数.
    设 \(w_1\), \(w_2\), \(\dots\), \(w_n\) 是数.
    则
    \begin{equation}
        \prod_{i = 1}^{n} {(x + w_i)}
        = x^n + \sum_{k = 1}^{n} {x^{n-k}
            \sum_{1 \leq \ell_1 < \dots < \ell_k \leq n}
            {w_{\ell_1} \dots w_{\ell_k}}
        }.
        \label{eq:Viète}
    \end{equation}
\end{theorem}

\begin{proof}
    作命题 \(P(n)\):
    对任何数 \(w_1\), \(w_2\), \(\dots\), \(w_n\),
    式~\eqref{eq:Viète} 是对的.

    \(P(1)\) 是对的.

    设 \(P(n)\) 是对的.
    我们由此证, \(P(n+1)\) 是对的.
    任取数 \(w_1\), \(\dots\), \(w_n\), \(w_{n+1}\),
    注意,
    \begin{align*}
        &
        \prod_{i = 1}^{n+1} {(x + w_i)}
        \\
        = {} &
        \left(\prod_{i = 1}^{n} {(x + w_i)}\right) (x + w_{n+1})
        \\
        = {} &
        \left(x^n + \sum_{k = 1}^{n} {x^{n-k}
                \sum_{1 \leq \ell_1 < \dots < \ell_k \leq n}
                {w_{\ell_1} \dots w_{\ell_k}}
        }\right) (x + w_{n+1})
        \\
        = {} &
        \hphantom{{} + {}}
        \left(x^n + \sum_{k = 1}^{n} {x^{n-k}
                \sum_{1 \leq \ell_1 < \dots < \ell_k \leq n}
                {w_{\ell_1} \dots w_{\ell_k}}
        }\right) x
        \\
        &
        {} +
        \left(x^n + \sum_{k = 1}^{n} {x^{n-k}
                \sum_{1 \leq \ell_1 < \dots < \ell_k \leq n}
                {w_{\ell_1} \dots w_{\ell_k}}
        }\right) w_{n+1}
        \\
        = {} &
        \hphantom{{} + {}}
        x^{n+1} + \sum_{k = 1}^{n} {x^{n+1-k}
            \sum_{1 \leq \ell_1 < \dots < \ell_k \leq n}
            {w_{\ell_1} \dots w_{\ell_k}}
        }
        \\
        &
        {} +
        x^n w_{n+1} + \sum_{k = 1}^{n} {x^{n-k}
            \sum_{1 \leq \ell_1 < \dots < \ell_k \leq n}
            {w_{\ell_1} \dots w_{\ell_k} w_{n+1}}
        }
        \\
        = {} &
        \hphantom{{} + {}}
        x^{n+1} + \sum_{u = 1}^{n} {x^{n+1-u}
            \sum_{1 \leq \ell_1 < \dots < \ell_u \leq n}
            {w_{\ell_1} \dots w_{\ell_u}}
        }
        \\
        &
        {} +
        x^n w_{n+1} + \sum_{u = 2}^{n+1} {x^{n-(u-1)}
            \sum_{1 \leq \ell_1 < \dots < \ell_{u-1} \leq n}
            {w_{\ell_1} \dots w_{\ell_{u-1}} w_{n+1}}
        }
        \\
        = {} &
        \hphantom{{} + {}}
        x^{n+1} + \sum_{u = 1}^{n} {x^{n+1-u}
            \sum_{1 \leq \ell_1 < \dots < \ell_u \leq n}
            {w_{\ell_1} \dots w_{\ell_u}}
        }
        \\
        &
        {} +
        x^n w_{n+1} + \sum_{u = 2}^{n+1} {x^{n+1-u}
            \sum_{\substack{
                    1 \leq \ell_1 < \dots < \ell_{u-1} < \ell_u \leq n+1 \\
                    \ell_u = n + 1
            }}
            {w_{\ell_1} \dots w_{\ell_{u-1}} w_{\ell_u}}
        }
        \\
        = {} &
        \hphantom{{} + {}}
        x^{n+1} + \sum_{u = 1}^{n} {x^{n+1-u}
            \sum_{1 \leq \ell_1 < \dots < \ell_u \leq n}
            {w_{\ell_1} \dots w_{\ell_u}}
        }
        \\
        &
        {}
        + x^{n+1-1}
        \sum_{\substack{
                1 \leq \ell_1 \leq n+1 \\
                \ell_1 = n + 1
        }}
        {w_{\ell_1} \dots w_{\ell_u}}
        + \sum_{u = 2}^{n} {x^{n+1-u}
            \sum_{\substack{
                    1 \leq \ell_1 < \dots < \ell_u \leq n+1 \\
                    \ell_u = n + 1
            }}
            {w_{\ell_1} \dots w_{\ell_u}}
        }
        \\
        &
        {}
        + x^{n+1-(n+1)} \sum_{\substack{
                1 \leq \ell_1 < \dots < \ell_{n+1} \leq n+1 \\
                \ell_{n+1} = n + 1
        }}
        {w_{\ell_1} \dots w_{\ell_u}}
        \\
        = {} &
        \hphantom{{} + {}}
        x^{n+1} + \sum_{u = 1}^{n} {x^{n+1-u}
            \sum_{\substack{
                    1 \leq \ell_1 < \dots < \ell_u \leq n+1 \\
                    \ell_u \leq n
            }}
            {w_{\ell_1} \dots w_{\ell_u}}
        }
        \\
        &
        {} +
        \sum_{u = 1}^{n} {x^{n+1-u}
            \sum_{\substack{
                    1 \leq \ell_1 < \dots < \ell_u \leq n+1 \\
                    \ell_u = n + 1
            }}
            {w_{\ell_1} \dots w_{\ell_{u-1}} w_{\ell_u}}
        }
        \\
        &
        {} +
        x^{n+1-(n+1)}
        \sum_{1 \leq \ell_1 < \dots < \ell_{n+1} \leq n+1}
        {w_{\ell_1} \dots w_{\ell_u}}
        \\
        = {} &
        \hphantom{{} + {}}
        x^{n+1} + \sum_{u = 1}^{n} {x^{n+1-u}
            \sum_{1 \leq \ell_1 < \dots < \ell_u \leq n+1}
            {w_{\ell_1} \dots w_{\ell_u}}
        }
        \\
        &
        {} +
        x^{n+1-(n+1)}
        \sum_{1 \leq \ell_1 < \dots < \ell_{n+1} \leq n+1}
        {w_{\ell_1} \dots w_{\ell_{n+1}}}
        \\
        = {} &
        x^{n+1} + \sum_{u = 1}^{n+1} {x^{n+1-u}
            \sum_{1 \leq \ell_1 < \dots < \ell_u \leq n+1}
            {w_{\ell_1} \dots w_{\ell_u}}
        }.
    \end{align*}
    所以, \(P(n+1)\) 是对的.
    由数学归纳法, 待证的命题成立.
\end{proof}

\begin{theorem}
    设 \(n\) 是正整数.
    设 \(x_1\), \(x_2\), \(\dots\), \(x_n\) 是数.
    设
    \begin{align*}
        f(x) = a_0 (x - x_1) (x - x_2) \dots (x - x_n).
    \end{align*}

    (1)
    若 \(f(x) = b_0 x^n + b_1 x^{n-1} + \dots + b_n\),
    则 \(b_0 = a_0\), 且若 \(1 \leq k \leq n\), 则
    \begin{align*}
        b_k = (-1)^k a_0 \sum_{1 \leq \ell_1 < \dots < \ell_k \leq n}
        {x_{\ell_1} \dots x_{\ell_k}}.
    \end{align*}

    (2)
    若 \(a_0 \neq 0\), 则 \(f(x)\) 是 \(n\)~次整式,
    且 \(x^n\)~项系数是 \(a_0\).

    (3)
    若数 \(c\) 是 \(x_1\), \(x_2\), \(\dots\), \(x_n\) 的一个,
    则 \(f(c) = 0\).
    反过来, 若 \(a_0 \neq 0\),
    且数 \(d\) 不等于 \(x_1\), \(x_2\), \(\dots\), \(x_n\),
    则 \(f(d) \neq 0\).
\end{theorem}

\begin{proof}
    (1)
    由上个定理, \(b_0 = a_0 \cdot 1 = a_0\),
    且若 \(1 \leq k \leq n\), 则
    \begin{align*}
        b_k
        = {} &
        a_0
        \sum_{1 \leq \ell_1 < \dots < \ell_k \leq n}
        {(-x_{\ell_1}) \dots (-x_{\ell_k})}
        \\
        = {} &
        a_0
        \sum_{1 \leq \ell_1 < \dots < \ell_k \leq n}
        {(-1)^k x_{\ell_1} \dots x_{\ell_k}}
        \\
        = {} &
        a_0\, (-1)^k
        \sum_{1 \leq \ell_1 < \dots < \ell_k \leq n}
        {x_{\ell_1} \dots x_{\ell_k}}
        \\
        = {} &
        (-1)^k a_0
        \sum_{1 \leq \ell_1 < \dots < \ell_k \leq n}
        {x_{\ell_1} \dots x_{\ell_k}}.
    \end{align*}

    (2)
    由 (1), \(f(x)\) 是次不高于 \(n\) 的整式,
    且 \(x^n\)~项系数 \(b_0 = a_0 \neq 0\).
    则 \(f(x)\) 是 \(n\)~次整式.

    (3)
    若数 \(c\) 是 \(x_1\), \(x_2\), \(\dots\), \(x_n\) 的一个,
    则 \(c - x_1\), \(c - x_2\), \(\dots\), \(c - x_n\) 的一个是 \(0\).
    注意, 任何数跟 \(0\) 的积是 \(0\).
    反过来, 若 \(a_0 \neq 0\),
    且数 \(d\) 不等于 \(x_1\), \(x_2\), \(\dots\), \(x_n\),
    则 \(a_0\), \(d - x_1\), \(d - x_2\), \(\dots\), \(d - x_n\)
    都非零.
    注意, 非零的数的积不是零.
\end{proof}

\section{结式的应用}

本节, 我们看结式的应用.

为方便, 我们仍作如下假定.

\AssumptionResultant*

设假定~\ref{standing-assumption-of-resultant} 成立.
回想, \(f\), \(g\) 的 Sylvester 阵是 \(m + n\)~级阵
\(S(f, g; n, m)\),
其中,
\begin{align*}
    [S(f, g; n, m)]_{i,j} =
    \begin{cases}
        a_{i-j},     & j \leq m; \\
        b_{i-(j-m)}, & j > m.
    \end{cases}
\end{align*}
\(f\), \(g\) 的结式是
\(f\), \(g\) 的 Sylvester 阵的行列式
\(\det {(S(f, g; n, m))}\).

在本节, 结式的以下的性质是重要的
(设假定~\ref{standing-assumption-of-resultant} 成立):

(1)
若存在数 \(c\), 使 \(f(c) = 0 = g(c)\),
则 \(\det {(S(f, g; n, m))} = 0\).

(2)
若 \(a_0 = 0 = b_0\),
则 \(\det {(S(f, g; n, m))} = 0\).

(3)
\(\det {(S(g, f; m, n))} = (-1)^{m n} \det {(S(f, g; n, m))}\).

(4)
若存在数 \(x_1\), \(x_2\), \(\dots\), \(x_n\),
使对任何数 \(x\),
\(f(x) = a_0 (x - x_1) \dots (x - x_n)\),
则 \(\det {(S(f, g; n, m))} = a_0^m g(x_1) \dots g(x_n)\);
若存在数 \(y_1\), \(y_2\), \(\dots\), \(y_m\),
使对任何数 \(x\),
\(g(x) = b_0 (x - y_1) \dots (x - y_m)\),
则 \(\det {(S(f, g; n, m))} = (-1)^{m n} b_0^n f(y_1) \dots f(y_m)\).

(5)
若 \(\det {(S(f, g; n, m))} = 0\),
且 \(a_0\), \(b_0\) 的一个非零
(它们当然可以都非零),
则存在数 \(c\), 使 \(f(c) = 0 = g(c)\).

若数 \(c\) 适合 \(f(c) = 0\),
我们说, \(c\) (或 \(x = c\)) 是 \(f(x) = 0\) 的一个解.
为方便, 我们也说, \(c\) 是 \(f\) (或 \(f(x)\)) 的一个\emph{根}.
于是, 通俗地, 我们可以说,
\(2\)~个整式 \(f\), \(g\) 有公根
(也就是, 存在数 \(c\), 使 \(c\) 既是 \(f\) 的根,
且 \(c\) 也是 \(g\) 的根),
相当于它们的结式是 \(0\)
(若 \(a_0\) 或 \(b_0\) 非零).

\begin{example}
    解方程组
    \begin{equation}
        \left\{ % tex-fmt: skip
        \begin{aligned}
            11x^2 - 2xy - 44y^2 - x + 26y - 32  & = 0, \\
            4x^2 - 18xy + 49y^2 - 4x + 9y - 118 & = 0.
        \end{aligned}
        \right. % tex-fmt: skip
        \label{eq:soe1}
    \end{equation}

    设 \(x = a\), \(y = b\) 是%
    方程组~\eqref{eq:soe1} 的一个解.
    则
    \begin{align*}
        &
        \begin{aligned}
            0
            = {} & 11a^2 - 2ab - 44b^2 - a + 26b - 32             \\
            = {} & 11 a^2 + (-2 b - 1) a + (-44 b^2 + 26 b - 32),
        \end{aligned}
        \\
        &
        \begin{aligned}
            0
            = {} & 4a^2 - 18ab + 49b^2 - 4a + 9b - 118           \\
            = {} & 4 a^2 + (-18 b - 4) a + (49 b^2 + 9 b - 118).
        \end{aligned}
    \end{align*}
    我们记
    \begin{align*}
        f(x; y) & = 11 x^2 + (-2 y - 1) x + (-44 y^2 + 26 y - 32), \\
        g(x; y) & = 4 x^2 + (-18 y - 4) x + (49 y^2 + 9 y - 118).
    \end{align*}
    则 \(f(a; b) = 0 = g(a; b)\).
    故 (关于 \(x\) 的) 整式
    \(f(x; b)\), \(g(x; b)\) 的结式是 \(0\),
    即
    \begin{align*}
        \begin{bmatrix}
            11            & 0             & 4            & 0            \\
            -2b-1         & 11            & -18b-4       & 4            \\
            -44b^2+26b-32 & -2b-1         & 49b^2+9b-118 & -18b-4       \\
            0             & -44b^2+26b-32 & 0            & 49b^2+9b-118 \\
        \end{bmatrix}
    \end{align*}
    的行列式是 \(0\).
    可以算出, 此阵的行列式是
    \(342\,125\, (b^2 - 1)\, (b^2 - 4)\).
    于是, 结式等于 \(0\),
    相当于 \(b = 1\) 或 \(b = -1\)
    或 \(b = 2\) 或 \(b = -2\).
    所以, 若 \(x = a\), \(y = b\) 是%
    方程组~\eqref{eq:soe1} 的一个解,
    则 \(b = 1\) 或 \(b = -1\)
    或 \(b = 2\) 或 \(b = -2\).

    另一方面,
    注意, \(f(x; b)\) 的 \(x^2\)~项系数 \(11 \neq 0\).
    所以, 若结式等于 \(0\)
    (也就是, \(b = 1\) 或 \(b = -1\)
    或 \(b = 2\) 或 \(b = -2\)),
    则存在数 \(a\), 使 \(f(a; b) = 0 = g(a; b)\).
    则 \(x = a\), \(y = b\)
    是方程组~\eqref{eq:soe1} 的解.
    为求 \(a\), 我们代方程组~\eqref{eq:soe1} 的
    \(b\) 以 \(1\), 得
    \begin{equation}
        \left\{ % tex-fmt: skip
        \begin{aligned}
            11a^2 - 2a - 44 - a + 26 - 32  & = 0, \\
            4a^2 - 18a + 49 - 4a + 9 - 118 & = 0;
        \end{aligned}
        \right. % tex-fmt: skip
        \label{eq:soe2}
    \end{equation}
    我们代 \(b\) 以 \(-1\), 得
    \begin{equation}
        \left\{ % tex-fmt: skip
        \begin{aligned}
            11a^2 + 2a - 44 - a - 26 - 32  & = 0, \\
            4a^2 + 18a + 49 - 4a - 9 - 118 & = 0;
        \end{aligned}
        \right. % tex-fmt: skip
        \label{eq:soe3}
    \end{equation}
    我们代 \(b\) 以 \(2\), 得
    \begin{equation}
        \left\{ % tex-fmt: skip
        \begin{aligned}
            11a^2 - 4a - 176 - a + 52 - 32   & = 0, \\
            4a^2 - 36a + 196 - 4a + 18 - 118 & = 0;
        \end{aligned}
        \right. % tex-fmt: skip
        \label{eq:soe4}
    \end{equation}
    我们代 \(b\) 以 \(-2\), 得
    \begin{equation}
        \left\{ % tex-fmt: skip
        \begin{aligned}
            11a^2 + 4a - 176 - a - 52 - 32   & = 0, \\
            4a^2 + 36a + 196 - 4a - 18 - 118 & = 0.
        \end{aligned}
        \right. % tex-fmt: skip
        \label{eq:soe5}
    \end{equation}
    方程组~\eqref{eq:soe2} 的解是 \(a = -2\);
    方程组~\eqref{eq:soe3} 的解是 \(a = 3\);
    方程组~\eqref{eq:soe4} 的解是 \(a = 4\);
    方程组~\eqref{eq:soe5} 的解是 \(a = -5\).
    于是, 方程组~\eqref{eq:soe1} 的解应是:
    \begin{equation*}
        \begin{alignedat}{7}
            {}            & x & {} = {} & {-2}, & \quad & y & {} = {} & 1; \\
            \text{或} \quad & x & {} = {} & 3,    & \quad & y & {} = {} & {-1};              \\
            \text{或} \quad & x & {} = {} & 4,    & \quad & y & {} = {} & 2;                 \\
            \text{或} \quad & x & {} = {} & {-5}, & \quad & y & {} = {} & {-2}.
        \end{alignedat}
    \end{equation*}
    最后, 不难验证, 这些确实都是解.
\end{example}

一般地, 我们能以结式解 \(2\)~元整式方程组
\begin{equation}
    \left\{ % tex-fmt: skip
    \begin{aligned}
        f(x; y) & = 0, \\
        g(x; y) & = 0.
    \end{aligned}
    \right. % tex-fmt: skip
    \label{eq:soe6}
\end{equation}
的解.
设我们已按 \(x\) 的次方从大到小地排 \(2\)~元整式的项, 即
\begin{align*}
    f(x; y) & = a_0 (y) x^n + a_1 (y) x^{n-1} + \dots + a_n (y), \\
    g(x; y) & = b_0 (y) x^m + b_1 (y) x^{m-1} + \dots + b_m (y),
\end{align*}
其中, \(m\), \(n\) 是非负整数,
\(m + n \geq 1\)
(注意, 若 \(m + n < 1\), % tex-fmt: skip
则 \(m = n = 0\);
也就是, \(f(x; y)\) 与 \(g(x; y)\) 不含 \(x\);
由此, 方程组~\eqref{eq:soe6}
其实是关于 \(y\) 的 \(1\)~元整式方程组), % tex-fmt: skip
且 \(a_0 (y)\), \(a_1 (y)\), \(\dots\), \(a_n (y)\),
\(b_0 (y)\), \(b_1 (y)\), \(\dots\), \(b_m (y)\)
是关于 \(y\) 的整式.
(注意, 我们可以要求 \(a_0 (y)\) 与 \(b_0 (y)\)
不是 \(0\) 整式.)
为方便, 我们约定:
若 \(k < 0\) 或 \(k > n\), 则整式 \(a_k (y) = 0\);
若 \(k < 0\) 或 \(k > m\), 则整式 \(b_k (y) = 0\).
由此, 我们写出 \(f\), \(g\) 的 Sylvester 阵
\(S(f, g; n, m)\).
它是 \(m + n\)~级阵, 且
\begin{align*}
    [S(f, g; n, m)]_{i,j} =
    \begin{cases}
        a_{i-j} (y),     & j \leq m; \\
        b_{i-(j-m)} (y), & j > m.
    \end{cases}
\end{align*}
则 \(f\), \(g\) 的结式 \(\det {(S(f, g; n, m))}\)
是关于 \(y\) 的整式.
那么, 若 \((x, y) = (x_0, y_0)\)
是方程组~\eqref{eq:soe6} 的解,
则 \(y_0\) 是 \(\det {(S(f, g; n, m))} = 0\) 的解;
也就是, 若 \(y_0\) 不是 \(\det {(S(f, g; n, m))} = 0\) 的解,
则对任何数 \(c\),
\(f(c; y_0) \neq 0\), 或 \(g(c; y_0) \neq 0\).
所以, 若 \(\det {(S(f, g; n, m))} = 0\) 无解,
则方程组~\eqref{eq:soe6} 也无解.
否则, 设 \(y_1\), \(\dots\), \(y_r\) 是
\(\det {(S(f, g; n, m))} = 0\) 的解.
代方程组~\eqref{eq:soe6} 的 \(y\) 以 \(y_i\)
(\(i = 1\), \(\dots\), \(r\)),
得关于 \(x\) 的 \(1\)~元整式方程组
\begin{equation*}
    \left\{ % tex-fmt: skip
    \begin{aligned}
        f(x; y_i) & = 0, \\
        g(x; y_i) & = 0.
    \end{aligned}
    \right. % tex-fmt: skip
\end{equation*}
解它, 得解 \(x_{i,1}\), \(\dots\), \(x_{i,k_i}\)
(注意, % tex-fmt: skip
若 \(a_0 (y_i) \neq 0\) 或 \(b_0 (y_i) \neq 0\),
则此关于 \(x\) 的 \(1\)~元整式方程组一定有解;
若 \(a_0 (y_i) = 0 = b_0 (y_i)\),
则它可能无解.) % tex-fmt: skip
则 \((x, y) = (x_{i,1}, y_i)\), \(\dots\), \((x_{i,k_i}, y_i)\)
是方程组~\eqref{eq:soe6} 的解,
\(i = 1\), \(\dots\), \(r\).

\begin{example}
    设数 \(x\), \(y\), \(t\) 适合
    \begin{equation}
        x = \frac{1 + 8t - t^2}{1 + t^2},
        \quad
        y = \frac{4t}{1 + t^2}.
        \label{eq:soe7}
    \end{equation}
    我们可消去 \(t\), 以得到 \(x\) 与 \(y\) 的关系:
    若 \((x, y, t) = (x_0, y_0, t_0)\) 适合式~\eqref{eq:soe7},
    则 \((x, y) = (x_0, y_0)\) 适合某关系 (不含 \(t\)),
    且反过来,
    若 \((x, y) = (x_0, y_0)\) 适合这个不含 \(t\) 的关系,
    则存在数 \(t_0\),
    使 \((x, y, t) = (x_0, y_0, t_0)\) 适合式~\eqref{eq:soe7}.

    设 \((x, y, t) = (x_0, y_0, t_0)\) 适合式~\eqref{eq:soe7}.
    去分母, 得
    \begin{align*}
        (1 + t_0^2) x_0 - (1 + 8t_0 - t_0^2) & = 0, \\
        (1 + t_0^2) y_0 - 4t_0               & = 0,
    \end{align*}
    即
    \begin{align*}
        (x_0 + 1) t_0^2 + (-8) t_0 + (x_0 - 1) & = 0, \\
        y_0 t_0^2 + (-4) t_0 + y_0             & = 0.
    \end{align*}
    记
    \begin{align*}
        f(t; x) & = (x + 1) t^2 + (-8) t + (x - 1), \\
        g(t; y) & = y t^2 + (-4) t + y.
    \end{align*}
    则 \(f(t_0; x_0)\) 与 \(g(t_0; y_0)\)
    等于 \(0\).
    反过来,
    若 \((x, y, t) = (x_0, y_0, t_0)\) 适合
    \(f(t_0; x_0) = 0 = g(t_0; y_0)\),
    则 \(t_0^2 + 1 \neq 0\).
    (反证法: 若 \(y_0 t_0^2 - 4 t_0 + y_0 = 0\), % tex-fmt: skip
    且 \(t_0^2 + 1 = 0\),
    则 \(4 t_0 = 0\).
    则 \(t_0 = 0\).
    可是, \(0^2 + 1 \neq 0\).) % tex-fmt: skip
    则 \((x, y, t) = (x_0, y_0, t_0)\) 适合式~\eqref{eq:soe7}.

    所以,
    若 \((x, y, t) = (x_0, y_0, t_0)\) 适合式~\eqref{eq:soe7},
    则 \(t_0\) 是 (关于 \(t\) 的) 整式
    \(f(t; x_0)\) 与 \(g(t; y_0)\)
    的公根.
    则 \(f(t; x_0)\) 与 \(g(t; y_0)\)
    的结式是 \(0\), 即
    \begin{align*}
        0 =
        \det {
            \begin{bmatrix}
                x_0+1 & 0     & y_0  & 0   \\
                -8    & x_0+1 & -4   & y_0 \\
                x_0-1 & -8    & y_0  & -4  \\
                0     & x_0-1 & 0    & y_0 \\
            \end{bmatrix}
        }
        = 4 (4 x_0^2 - 16 x_0 y_0 + 17 y_0^2 - 4).
    \end{align*}
    记
    \begin{align*}
        R(x, y) = 4 (4 x^2 - 16 x y + 17 y^2 - 4).
    \end{align*}
    则 \(R(x, y) = 0\).
    反过来,
    若 \((x, y) = (x_0, y_0)\) 适合 \(R(x, y) = 0\),
    则 \(x_0 + 1 = 0 = y_0\),
    或当 \(x_0 + 1 \neq 0\) 或 \(y_0 \neq 0\) 时,
    存在数 \(t_0\),
    使 \(f(t_0; x_0) = 0 = g(t_0; y_0)\).
    所以,
    若 \((x, y) = (x_0, y_0)\) 适合 \(R(x_0, y_0) = 0\),
    则 \((x_0, y_0) = (-1, 0)\),
    或当 \((x_0, y_0) \neq (-1, 0)\) 时,
    存在数 \(t_0\),
    使 \((x, y, t) = (x_0, y_0, t_0)\) 适合式~\eqref{eq:soe7}.
    我们可以验证, 不存在数 \(t_0\),
    使 \((x, y, t) = (-1, 0, t_0)\) 适合式~\eqref{eq:soe7}.
    (反证法: 若 \(0 = 4t_0/(1 + t_0^2)\), % tex-fmt: skip
    则 \(t_0 = 0\).
    可是, \(-1 \neq (1 + 8 \cdot 0 - 0^2) % tex-fmt: skip
    / (1 + 0^2)\).) % tex-fmt: skip

    综上, \(x\), \(y\) 的关系是
    (我们以消去律消去 \(4\))
    \begin{equation*}
        4 x^2 - 16 x y + 17 y^2 - 4 = 0,
        \quad
        (x, y) \neq (-1, 0).
        \qefhere
    \end{equation*}
\end{example}

一般地, 我们能以结式从 \(3\)~元 \(x\), \(y\), \(t\)
适合的关系
\begin{equation}
    x = \frac{q_1 (t)}{p_1 (t)},
    \quad
    y = \frac{q_2 (t)}{p_2 (t)},
    \label{eq:soe8}
\end{equation}
其中, \(q_1\), \(q_2\), \(p_1\), \(p_2\)
是关于 \(t\) 的整式,
且 \(p_1\), \(p_2\) 不是 \(0\) 整式 (分母不为 \(0\)),
且 \(q_1\), \(q_2\) 不是 \(0\) 整式
(若, 比如, \(q_1 = 0\), 则 \(x = 0\);
这是平凡的),
消去 \(t\), 以得到
\(2\)~元 \(x\) 与 \(y\) 适合的关系.
注意, 我们可以设 \(p_1\), \(q_1\) 没有公根,
且 \(p_2\), \(q_2\) 没有公根
(因为, 若, 比如, % tex-fmt: skip
\(p_1 (c) = 0 = q_1 (c)\),
则存在非零的整式 \(r_{1,1}\), \(s_{1,1}\),
使 \(p_1 (x) = (x - c) r_{1,1} (x)\),
\(q_1 (x) = (x - c) s_{1,1} (x)\),
且 \(r_{1,1}\) 的次低于 \(p_1\) 的次,
且 \(s_{1,1}\) 的次低于 \(q_1\) 的次;
则 \(q_1 (t) / p_1 (t) = r_{1,1} (t) / s_{1,1} (t)\);
若 \(r_{1,1}\), \(s_{1,1}\) 没有公根, 则那是好的;
若 \(r_{1,1}\), \(s_{1,1}\) 有公根,
则类似地, 存在非零的整式 \(r_{1,2}\), \(s_{1,2}\),
使 \(r_{1,2}\) 的次低于 \(r_{1,1}\) 的次,
且 \(s_{1,2}\) 的次低于 \(s_{1,1}\) 的次,
且 \(q_1 (t) / p_1 (t)
    = r_{1,1} (t) / s_{1,1} (t)
= r_{1,2} (t) / s_{1,2} (t)\);
非零的整式的次不能无限地变少,
所以这个过程能结束;
最后, 注意, 即使这改变了 \(t\) 的允许的取值,
我们可另外讨论此事). % tex-fmt: skip
去分母, 且移项, 得
\begin{equation}
    x p_1 (t) - q_1 (t) = 0,
    \quad
    y p_2 (t) - q_2 (t) = 0.
    \label{eq:soe9}
\end{equation}
注意, 若 \((x, y, t) = (x_0, y_0, t_0)\)
适合式~\eqref{eq:soe8},
则去分母, 且移项, 我们知道
\((x, y, t) = (x_0, y_0, t_0)\)
也适合式~\eqref{eq:soe9};
反过来, 若 \((x, y, t) = (x_0, y_0, t_0)\)
适合式~\eqref{eq:soe9},
则注意, \(p_1 (t_0) \neq 0\), 且 \(p_2 (t_0) \neq 0\)
(因为, 若, 比如, \(p_1 (t_0) = 0\), % tex-fmt: skip
则 \(q_1 (t_0) = x_0 p_1 (t_0) = 0\);
不过, 我们已设, \(p_1\), \(q_1\) 没有公根,
故这不是可能的). % tex-fmt: skip
则作除法, 且移项, 我们知道,
\((x, y, t) = (x_0, y_0, t_0)\)
也适合式~\eqref{eq:soe8}.
然后, 设
\begin{align*}
    f(t; x) & = x p_1 (t) - q_1 (t)
    = a_0 (x) t^n + a_1 (x) t^{n-1} + \dots + a_n (x),
    \\
    g(t; y) & = y p_2 (t) - q_2 (t)
    = b_0 (y) t^m + b_1 (y) t^{m-1} + \dots + b_m (y),
\end{align*}
其中, \(m\), \(n\) 是非负整数,
\(m + n \geq 1\)
(注意, 若 \(m + n < 1\), % tex-fmt: skip
则 \(m = n = 0\);
也就是, \(f(t; x)\) 与 \(g(t; y)\) 不含 \(t\);
这是平凡的), % tex-fmt: skip
\(a_0 (x)\), \(a_1 (x)\), \(\dots\), \(a_n (x)\)
是关于 \(x\) 的整式,
且 \(b_0 (y)\), \(b_1 (y)\), \(\dots\), \(b_m (y)\)
是关于 \(y\) 的整式.
(注意, 我们可以要求 \(a_0 (x)\) 与 \(b_0 (y)\)
不是 \(0\) 整式.)
为方便, 我们约定:
若 \(k < 0\) 或 \(k > n\), 则整式 \(a_k (x) = 0\);
若 \(k < 0\) 或 \(k > m\), 则整式 \(b_k (y) = 0\).
由此, 我们写出 \(f\), \(g\) 的 Sylvester 阵
\(S(f, g; n, m)\).
它是 \(m + n\)~级阵, 且
\begin{align*}
    [S(f, g; n, m)]_{i,j} =
    \begin{cases}
        a_{i-j} (x),     & j \leq m; \\
        b_{i-(j-m)} (y), & j > m.
    \end{cases}
\end{align*}
则 \(f\), \(g\) 的结式 \(\det {(S(f, g; n, m))}\)
是关于 \(x\), \(y\) 的整式.
那么, 若 \((x, y, t) = (x_0, y_0, t_0)\)
适合式~\eqref{eq:soe8},
则它适合式~\eqref{eq:soe9},
故 \((x, y) = (x_0, y_0)\)
适合 \(\det {(S(f, g; n, m))} = 0\);
也就是, 若 \((x, y) = (x_0, y_0)\) 适合
\(\det {(S(f, g; n, m))} \neq 0\),
则对任何数 \(t_0\),
\((x, y, t) = (x_0, y_0, t_0)\)
不适合式~\eqref{eq:soe8}.
反过来, 若 \((x, y) = (x_0, y_0)\) 适合
\(\det {(S(f, g; n, m))} = 0\),
则 \(a_0 (x_0) = 0 = b_0 (y_0)\),
或当 \(a_0 (x_0) \neq 0\) 或 \(b_0 (y_0) \neq 0\) 时,
存在数 \(t_0\), 使
\((x, y, t) = (x_0, y_0, t_0)\)
适合式~\eqref{eq:soe9},
故它适合式~\eqref{eq:soe8}.
当然, 对适合 \(a_0 (x_0) = 0 = b_0 (y_0)\) 的
\((x_0, y_0)\),
我们可另外讨论, 是否存在数 \(t_0\),
使 \((x, y, t) = (x_0, y_0, t_0)\)
适合式~\eqref{eq:soe8}.

\begin{example}
    设数 \(c\) 适合 \(c^2 + c = 1\).
    求 \(c^4 - c^2 + 2c + 1\).

    一个方法或许是先确定 \(c\),
    再由此求 \(c^2\), \(c^4\),
    再求 \(c^4 - c^2 + 2c + 1\).

    这里, 我们以结式解此事.
    为方便, 我们记
    \(f(x) = 1 x^2 + 1 x - 1\),
    \(v(x) = x^4 - x^2 + 2x + 1\).
    则问题相当于, 若 \(f(c) = 0\), 则 \(v(c)\) 是何.
    我们记 \(v(c) = y\), 且记
    \(g(x) = y - v(x)
    = -1 x^4 + 0 x^3 + 1 x^2 + (-2) x + (y - 1)\).
    则 \(f(c) = 0 = g(c)\).
    所以, \(f(x)\), \(g(x)\) 的结式是 \(0\),
    即
    \begin{align*}
        0 =
        \det {
            \begin{bmatrix}
                1  & 0  & 0  & 0  & -1  & 0   \\
                1  & 1  & 0  & 0  & 0   & -1  \\
                -1 & 1  & 1  & 0  & 1   & 0   \\
                0  & -1 & 1  & 1  & -2  & 1   \\
                0  & 0  & -1 & 1  & y-1 & -2  \\
                0  & 0  & 0  & -1 & 0   & y-1 \\
            \end{bmatrix}
        }
        = (y - 2)^2.
    \end{align*}
    故 \(v(c) = y = 2\).
\end{example}

一般地, 我们能以结式求值.
设 \(n \geq 1\), 且 \(a_0 \neq 0\).
设整式 \(f(x) = a_0 x^n + a_1 x^{n-1} + \dots + a_n\).
设 \(u(x)\), \(v(x)\) 是整式.
设数 \(c\) 适合 \(f(c) = 0\),
且 \(u(c) \neq 0\).
我们求 \(v(c) / u(c)\).
为此, 我们设 \(y = v(c) / u(c)\),
并作整式
\(g(x) = y u(x) - v(x)
= b_0 (y) x^m + b_1 (y) x^{m-1} + \dots + b_m (y)\).
则 \(g(c) = (v(c) / u(c)) u(c) - v(c) = 0\).
为方便, 我们约定:
若 \(k < 0\) 或 \(k > n\), 则 \(a_k = 0\);
若 \(k < 0\) 或 \(k > m\), 则整式 \(b_k (y) = 0\).
由此, 我们写出 \(f\), \(g\) 的 Sylvester 阵
\(S(f, g; n, m)\).
它是 \(m + n\)~级阵, 且
\begin{align*}
    [S(f, g; n, m)]_{i,j} =
    \begin{cases}
        a_{i-j},         & j \leq m; \\
        b_{i-(j-m)} (y), & j > m.
    \end{cases}
\end{align*}
则 \(f\), \(g\) 的结式 \(\det {(S(f, g; n, m))}\)
是关于 \(y\) 的整式.
因为 \(f(c) = 0 = g(c)\),
\(\det {(S(f, g; n, m))} = 0\).
这是一个关于 \(y\) 的方程.
解它, 得 \(y = v(c) / u(c)\) 的所有的可能的取值.

\begin{definition}
    设 \(f(x)\) 是整式.
    设 \(c\) 是数.

    (1)
    若 \(f(c) = 0\), 则 \(c\) 是 \(f(x)\) 的一个\emph{根}.

    (2)
    若 \(f(c) = 0\), 且 \(\mathrm{D} f(c) = 0\),
    则 \(c\) 是 \(f(x)\) 的一个\emph{重根}.
\end{definition}

\begin{example}
    设 \(a\), \(b\), \(c\) 是实数, 且 \(a \neq 0\).
    设 \(f(x) = a x^2 + b x + c\);
    注意, 这是 \(2\)~次函数.
    回想, 在中学, 我们说,
    若 \(1\)~元 \(2\)~次方程
    \(a x^2 + b x + c = 0\) 的%
    判别式 \(b^2 - 4 a c\) 等于 \(0\),
    则 \(a x^2 + b x + c = 0\) 有 ``重根''
    \begin{align*}
        x_1 = x_2 = -\frac{b}{2a}.
    \end{align*}
    那么, 这跟我们定义的重根有何关系?
    我们记 \(r = -b/(2a)\).
    则 \(b = -2a r\),
    且 \(c = b^2 / (4a) = a r^2\).
    则 \(f(x) = a x^2 - 2 a r x + a r^2
    = a (x^2 - 2 r x + r^2)\).
    首先, \(r\) 是 \(f(x)\) 的根:
    \(f(r) = a (r^2 - 2 r r + r^2) = 0\);
    然后, \(\mathrm{D} f(x) = a (2 x - 2 r)\),
    故 \(\mathrm{D} f(r) = a (2 r - 2 r) = 0\).
    故, 由定义, \(r\) 是 \(f(x)\) 的重根.
\end{example}

\begin{example}
    每个数都是 \(0\) 整式的根.
    注意, \(0\) 的 derivative 是 \(0\).
    于是, 每个数都是 \(0\) 整式的重根.

    设 \(f(x) = c\) (\(c \neq 0\)) 是 \(0\)~次整式.
    则 \(f(x)\) 没有根.
    则 \(f(x)\) 没有重根.

    设 \(g(x) = a x + b\) (\(a \neq 0\)) 是 \(1\)~次整式.
    则 \(-b/a\) 是 \(g(x)\) 的根,
    因为 \(g(-b/a) = a (-b/a) + b = -b + b = 0\).
    不过, 因为 \(\mathrm{D} g(x) = a \neq 0\),
    \(g(x)\) 没有重根.
\end{example}

此定理或许解释了为何我们叫重根为 ``重'' 根------重根%
真地重出现了.

\begin{theorem}
    设 \(f(x)\) 是整式.

    (1)
    若数 \(c\) 是 \(f(x)\) 的重根,
    则存在整式 \(s(x)\),
    使 \(f(x) = (x - c)^2 s(x)\).

    (2)
    若数 \(d\) 与整式 \(u(x)\) 适合,
    \(f(x) = (x - d)^2 u(x)\),
    则 \(d\) 是 \(f(x)\) 的重根.
\end{theorem}

\begin{proof}
    (1)
    因为 \(c\) 是 \(f(x)\) 的重根,
    故 \(c\) 是 \(f(x)\) 的根.
    则存在整式 \(q(x)\), 使 \(f(x) = (x - c) q(x)\).
    注意, \(0 = \mathrm{D} f(c) = q(c)\).
    则存在整式 \(s(x)\), 使 \(q(x) = (x - c) s(x)\).
    故 \(f(x) = (x - c)^2 s(x)\).

    (2)
    记 \(v(x) = (x - d) u(x)\).
    则 \(f(x) = (x - d) v(x)\).
    则 \(f(d) = 0\),
    且 \(\mathrm{D} f(d) = v(d) = 0\).
    故 \(d\) 是 \(f(x)\) 的重根.
\end{proof}

我们讨论整式 \(f(x)\) 与它的 derivative \(\mathrm{D} f(x)\)
的 Sylvester 阵与结式.
设 \(f(x) = a_0 x^n + a_1 x^{n-1} + \dots + a_{n-1} x + a_n\).
则 \(\mathrm{D} f(x)
= n a_0 x^{n-1} + (n-1) a_1 x^{n-2} + \dots + a_{n-1}\).
为方便, 我们约定:
若 \(k < 0\) 或 \(k > n\), 则 \(a_k = 0\).
记 \(b_k = (n - k) a_k\).
则 \(\mathrm{D} f(x)
= b_0 x^{n-1} + b_1 x^{n-2} + \dots + b_{n-1}\),
且若 \(k < 0\) 或 \(k > n-1\), 则 \(b_k = 0\).
注意, \((n - 1) + n \geq 1\), 相当于 \(n \geq 1\).
所以, 若 \(n \geq 1\), 我们能写出 \(f\), \(\mathrm{D} f\)
的 Sylvester 阵 \(S(f, \mathrm{D} f; n, n-1)\):
它是 \(2n-1\)~级阵, 且
(注意, \(b_{i-(j-(n-1))} = (1 - (i-j)) a_{i-j+n-1}\))
\begin{equation}
    [S(f, \mathrm{D} f; n, n-1)]_{i,j} =
    \begin{cases}
        a_{i-j},
        & j \leq n-1; \\
        (1 - (i - j)) a_{i-j+n-1},
        & j > n-1.
    \end{cases}
    \label{eq:soe10}
\end{equation}
注意, \([S(f, \mathrm{D} f; n, n-1)]_{1,1} = a_0\),
\([S(f, \mathrm{D} f; n, n-1)]_{1,n} = n a_0\),
且若 \(j \neq 1\) 且 \(j \neq n\),
则 \([S(f, \mathrm{D} f; n, n-1)]_{1,n} = 0\).
由此, 我们考虑 \(2n-1\)~级阵 \(T(f; n)\), 其中,
\begin{equation}
    [T(f; n)]_{i,j} =
    \begin{cases}
        1,
        & \text{\(i = 1\), 且 \(j = 1\)}; \\
        n,
        & \text{\(i = 1\), 且 \(j = n\)}; \\
        0,
        & \text{\(i = 1\), \(j \neq 1\), 且 \(j \neq n\)}; \\
        a_{i-j},
        & \text{\(i \neq 1\), 且 \(j \leq n-1\)}; \\
        (1 - (i - j)) a_{i-j+n-1},
        & \text{\(i \neq 1\), 且 \(j \geq n\)}.
    \end{cases}
    \label{eq:soe11}
\end{equation}
不难看出, 以 \(a_0\) 乘 \(T(f; n)\) 的行~\(1\),
得 \(S(f, \mathrm{D} f; n, n-1)\).
则
\begin{align*}
    \det {(S(f, \mathrm{D} f; n, n-1))}
    = a_0 \det {(T(f; n))}.
\end{align*}
因为 \(T(f; n)\) 的元全是关于
\(a_0\), \(a_1\), \(\dots\), \(a_n\) 的整式,
\(\det {(T(f; n))}\) 也是关于
\(a_0\), \(a_1\), \(\dots\), \(a_n\) 的整式.
注意, \(a_0\) 是 \(f\) 的 \(x^n\) 项系数;
注意, 我们可以设 \(a_0 \neq 0\)
(若 \(a_0 = 0\), % tex-fmt: skip
且 \(f(x)\) 不是 \(0\) 整式,
我们可以写 \(f(x)\)
为 \(a'_0 x^m + a'_1 x^{m-1} + \dots + a'_m\),
其中, \(m < n\), 且 \(a'_0 \neq 0\);
若 \(f(x)\) 是 \(0\) 整式, 则这是平凡的). % tex-fmt: skip
于是, 若 \(f\) 有重根,
则 \(f\) 与 \(\mathrm{D} f\) 有公根,
故 \(f\) 与 \(\mathrm{D} f\) 的结式为 \(0\);
由此, \(\det {(T(f; n))} = 0\).
另一方面, 若 \(\det {(T(f; n))} = 0\),
则 \(f\) 与 \(\mathrm{D} f\) 的结式为 \(0\);
因为 \(a_0 \neq 0\),
\(f\) 与 \(\mathrm{D} f\) 有公根;
由此, \(f\) 有重根.

设
\begin{align*}
    f(x) = a_0 (x - x_1) (x - x_2) \dots (x - x_n).
\end{align*}
我们具体地写 \(\det {(S(f, \mathrm{D} f; n, n-1))}\)
(与 \(\det {(T(f; n))}\)).
由结式的性质~(4),
\begin{align*}
    \det {(S(f, \mathrm{D} f; n, n-1))}
    = a_0^{n-1}
    \mathrm{D} f(x_1) \mathrm{D} f(x_2)
    \dots \mathrm{D} f(x_n).
\end{align*}
回想,
\begin{align*}
    \mathrm{D} f(x_i) = q_i (x_i)
    = {} &
    a_0 \prod_{\substack{1 \leq k \leq n \\ k \neq i}} {(x_i - x_k)}.
\end{align*}
则
\begin{align*}
    \det {(S(f, \mathrm{D} f; n, n-1))}
    = {} & a_0^{n-1} a_0^n
    \prod_{1 \leq i \leq n}
    \prod_{\substack{1 \leq k \leq n \\ k \neq i}} {(x_i - x_k)}
    \\
    = {} &
    a_0^{2n-1}
    \prod_{1 \leq k \neq i \leq n} {(x_i - x_k)}
    \\
    = {} &
    a_0^{2n-1}
    \prod_{1 \leq k < i \leq n} {(x_i - x_k)}
    \cdot \prod_{1 \leq i < k \leq n} {(x_i - x_k)}
    \\
    = {} &
    a_0^{2n-1}
    \prod_{1 \leq i < k \leq n} {(x_k - x_i)}
    \cdot \prod_{1 \leq i < k \leq n} {((-1) (x_k - x_i))}
    \\
    = {} &
    a_0^{2n-1}
    \prod_{1 \leq i < k \leq n} {((-1) (x_k - x_i)^2)}
    \\
    = {} &
    a_0^{2n-1}
    \prod_{1 \leq i < k \leq n} {(-1)}
    \cdot \prod_{1 \leq i < k \leq n} {(x_k - x_i)^2}
    \\
    = {} &
    a_0^{2n-1} (-1)^{(n-1)+(n-2)+\dots+1}
    \prod_{1 \leq i < k \leq n} {(x_k - x_i)^2}
    \\
    = {} &
    a_0
    \cdot (-1)^{n(n-1)/2} a_0^{2n-2}
    \prod_{1 \leq i < k \leq n} {(x_k - x_i)^2}.
\end{align*}
则
\begin{align*}
    \det {(T(f; n))}
    = (-1)^{n(n-1)/2} a_0^{2n-2}
    \prod_{1 \leq i < k \leq n} {(x_k - x_i)^2}.
\end{align*}

现在, 我们可以定义整式的判别式.

\begin{definition}
    设 \(f(x)\) 是整式.

    (1)
    若 \(f(x)\) 是 \(0\) 整式,
    则我们定义 \(f(x)\) 的\emph{判别式}为 \(0\).

    (2)
    若 \(f(x)\) 是非零常数 \(c\),
    则我们定义 \(f(x)\) 的\emph{判别式}为 \(c^{-2}\).

    (3)
    在别的情形, 我们可写
    \(f(x) = a_0 x^n + a_1 x^{n-1} + \dots + a_n\)
    (\(a_0 \neq 0\), 且 \(n \geq 1\)).
    我们定义 \(f(x)\) 的\emph{判别式}为
    \((-1)^{n (n-1)/2} \det {(T(f; n))}\),
    其中, \(T(f; n)\) 是 \(2n-1\)~级阵,
    且它的元由式~\eqref{eq:soe11} 确定.
\end{definition}

由前面的计算与讨论, 下面的结论是显然的.

\begin{theorem}
    设 \(n \geq 1\), 且 \(a_0 \neq 0\).
    设整式
    \begin{align*}
        f(x)
        = a_0 x^n + a_1 x^{n-1} + \dots + a_n
        = a_0 (x - x_1) (x - x_2) \dots (x - x_n).
    \end{align*}

    (1)
    \(f(x)\) 的判别式
    \begin{align*}
        (-1)^{n (n-1)/2} \det {(T(f; n))}
        = a_0^{2n-2}
        \prod_{1 \leq i < k \leq n} {(x_k - x_i)^2},
    \end{align*}
    其中, \(T(f; n)\) 是 \(2n-1\)~级阵,
    且它的元由式~\eqref{eq:soe11} 确定.

    \(f(x)\) 的判别式%
    是关于 \(a_0\), \(a_1\), \(\dots\), \(a_n\) 的整式.

    由式~\eqref{eq:soe10} 确定
    \(f\) 与 \(f\) 的 derivative \(\mathrm{D} f\)
    的 (\(2n-1\)~级) Sylvester 阵
    \(S(f, \mathrm{D} f; n, n-1)\).
    则 \(f\), \(\mathrm{D} f\) 的结式
    \begin{align*}
        \det {(S(f, \mathrm{D} f; n, n-1))}
        = (-1)^{n (n-1)/2} a_0
        \cdot (-1)^{n (n-1)/2} \det {(T(f; n))}.
    \end{align*}

    (2)
    若 \(f(x)\) 有重根,
    则 \(f(x)\) 的判别式为 \(0\);
    反过来, 若 \(f(x)\) 的判别式为 \(0\),
    则 \(f(x)\) 有重根.

    (3)
    若 \(a_0\) 是实数,
    且 \(x_1\), \(x_2\), \(\dots\), \(x_n\) 是实数,
    则 \(f(x)\) 的判别式是非负数.
\end{theorem}

\begin{example}
    设 \(a\), \(b\), \(c\) 是实数, 且 \(a \neq 0\).
    设 \(f(x) = a x^2 + b x + c\);
    注意, 这是 \(2\)~次函数.
    回想, 在中学, 我们说,
    \(1\)~元 \(2\)~次方程
    \(a x^2 + b x + c = 0\) 的
    ``判别式'' 是 \(b^2 - 4 a c\).
    那么, 这跟我们定义的判别式有何关系?
    为此, 我们写出
    \begin{align*}
        T(f; 2) =
        \begin{bmatrix}
            1 & 2 & 0   \\
            b & b & 2 a \\
            c & 0 & b   \\
        \end{bmatrix}.
    \end{align*}
    则 \(f(x)\) 的判别式
    \begin{align*}
        (-1)^{2 (2-1)/2} \det {
            \begin{bmatrix}
                1 & 2 & 0   \\
                b & b & 2 a \\
                c & 0 & b   \\
            \end{bmatrix}
        }
        = -(1 b^2 - 2 (b^2 - 2ac) + 0)
        = b^2 - 4ac.
    \end{align*}
    则 \(2\)~个判别式是相同的.
\end{example}

\begin{example}
    设 \(n \geq 2\).
    设 \(a_0 \neq 0\).
    设 \(f(x) = a_0 x^n + a_{n-1} x + a_n\).
    我们计算 \(f(x)\) 的判别式.

    为方便, 我们约定,
    若 \(k \neq 0\), \(n-1\), \(n\),
    则 \(a_k = 0\).
    回想, \(2n-1\)~级阵 \(T(f; n)\) 的元
    \begin{align*}
        [T(f; n)]_{i,j} =
        \begin{cases}
            1,
            & \text{\(i = 1\), 且 \(j = 1\)}; \\
            n,
            & \text{\(i = 1\), 且 \(j = n\)}; \\
            0,
            & \text{\(i = 1\), \(j \neq 1\), 且 \(j \neq n\)}; \\
            a_{i-j},
            & \text{\(i \neq 1\), 且 \(j \leq n-1\)}; \\
            (1 - (i - j)) a_{i-j+n-1},
            & \text{\(i \neq 1\), 且 \(j \geq n\)}.
        \end{cases}
    \end{align*}
    通俗地, 比如, 当 \(n = 5\) 时,
    \begin{align*}
        T(f; 5) =
        \begin{bmatrix}
            1   & 0   & 0   & 0   & 5   & 0     & 0     & 0     & 0     \\
            0   & a_0 & 0   & 0   & 0   & 5 a_0 & 0     & 0     & 0     \\
            0   & 0   & a_0 & 0   & 0   & 0     & 5 a_0 & 0     & 0     \\
            0   & 0   & 0   & a_0 & 0   & 0     & 0     & 5 a_0 & 0     \\
            a_4 & 0   & 0   & 0   & a_4 & 0     & 0     & 0     & 5 a_0 \\
            a_5 & a_4 & 0   & 0   & 0   & a_4   & 0     & 0     & 0     \\
            0   & a_5 & a_4 & 0   & 0   & 0     & a_4   & 0     & 0     \\
            0   & 0   & a_5 & a_4 & 0   & 0     & 0     & a_4   & 0     \\
            0   & 0   & 0   & a_5 & 0   & 0     & 0     & 0     & a_4   \\
        \end{bmatrix}.
    \end{align*}
    计算判别式前, 我们考虑, \(T(f; n)\) 的何元一定是 \(0\).
    若 \(i = 1\), 则当 \(j \neq 1\), \(n\) 时,
    \([T(f; n)]_{i,j} = 0\).
    若 \(i > 1\), 且 \(j \leq n - 1\),
    则当 \(i - j \neq 0\), \(n-1\), \(n\) 时,
    \([T(f; n)]_{i,j} = 0\),
    即当 \(j \leq n - 1\),
    且 \(i \neq j\), \(j + n - 1\), \(j + n\) 时,
    \([T(f; n)]_{i,j} = 0\).
    若 \(i > 1\), 且 \(j \geq n\),
    则当 \(i - j = 1\),
    或当 \(i - j + n - 1 \neq 0\), \(n-1\), \(n\) 时,
    \([T(f; n)]_{i,j} = 0\),
    即当 \(j \geq n\),
    且 \(i \neq j - n + 1\), \(j\) 时,
    \([T(f; n)]_{i,j} = 0\).
    综上,
    若 \(j \leq n-1\),
    则当 \(i \neq j\), \(j+n-1\), \(j+n\) 时,
    \([T(f; n)]_{i,j} = 0\);
    若 \(j \geq n\),
    则当 \(i \neq j-n+1\), \(j\) 时,
    \([T(f; n)]_{i,j} = 0\).

    注意, 若 \(i \leq n-1\), 且 \(j \leq n-1\),
    则 \([T(f; n)]_{i,j+n-1} = n [T(f; n)]_{i,j}\).
    注意, \(j + n > j + n - 1 \geq n > i\),
    故若 \(i \neq j\), 则 \([T(f; n)]_{i,j} = 0\);
    此时, \(j + n - 1 \geq n\),
    且 \(i \neq (j + n - 1) - n + 1\),
    即 \([T(f; n)]_{i,j+n-1} = 0\).
    若 \(i = j = 1\),
    则 \([T(f; n)]_{i,j} = 1\),
    且 \([T(f; n)]_{i,j+n-1} = n\);
    若 \(1 < i = j \leq n - 1\),
    则 \([T(f; n)]_{i,j} = a_0\),
    且 \([T(f; n)]_{i,j+n-1} = n a_0\).

    对 \(T(f; n)\),
    我们加列~\(c\) 的 \(-n\)~倍于列~\(c+n-1\),
    \(c = 1\), \(\dots\), \(n-1\),
    得阵 \(U\).
    则 \(U\) 的前 \(n-1\)~列与列~\(2n-1\)
    分别是 \(T(f; n)\) 的前 \(n-1\)~列与列~\(2n-1\),
    且若 \(n \leq j \leq 2n-2\),
    则当 \(i \leq n - 1\) 时, \([U]_{i,j} = 0\),
    且当 \(i = j\) 时,
    \([U]_{i,j} \allowbreak
        = [T(f; n)]_{i,j+n-1} - n [T(f; n)]_{i,j}
    = (1 - n) a_{n-1}\),
    且当 \(i = j + 1\) 时,
    \([U]_{i,j} \allowbreak
        = [T(f; n)]_{i,j+n-1} - n [T(f; n)]_{i,j}
    = -n a_n\),
    且当 \(n - 1 < i < j\) 或 \(i > j + 1\) 时,
    \([U]_{i,j} = 0\).
    综上,
    \begin{align*}
        [U]_{i,j} =
        \begin{cases}
            1,             & i = j = 1;                           \\
            a_0,           & 2 \leq i = j \leq n - 1;             \\
            a_{n-1},       & 2 \leq i - (n - 1) = j \leq n - 1;   \\
            a_n,           & 2 \leq i - n = j \leq n - 1;         \\
            (1-n) a_{n-1}, & n \leq i = j \leq 2n - 2;            \\
            -n a_n,        & n \leq i - 1 = j \leq 2n - 2;        \\
            n a_0,         & \text{\(i = n\), 且 \(j = 2n - 1\)}; \\
            a_{n-1},       & i = j = n;                           \\
            0,             & \text{别的情形}.
        \end{cases}
    \end{align*}
    通俗地, 比如, 当 \(n = 5\) 时,
    \begin{align*}
        U =
        \begin{bmatrix}
            1   & 0   & 0   & 0   & 0     & 0     & 0     & 0     & 0    \\
            0   & a_0 & 0   & 0   & 0     & 0     & 0     & 0     & 0    \\
            0   & 0   & a_0 & 0   & 0     & 0     & 0     & 0     & 0    \\
            0   & 0   & 0   & a_0 & 0     & 0     & 0     & 0     & 0    \\
            a_4 & 0   & 0   & 0   & -4a_4 & 0     & 0     & 0     & 5a_0 \\
            a_5 & a_4 & 0   & 0   & -5a_5 & -4a_4 & 0     & 0     & 0    \\
            0   & a_5 & a_4 & 0   & 0     & -5a_5 & -4a_4 & 0     & 0    \\
            0   & 0   & a_5 & a_4 & 0     & 0     & -5a_5 & -4a_4 & 0    \\
            0   & 0   & 0   & a_5 & 0     & 0     & 0     & -5a_5 & a_4  \\
        \end{bmatrix}.
    \end{align*}
    注意,
    \(\det {(T(f; n))} = \det {(U)}\).
    注意, \(U\)~的列~\(2n-1\) 有 \(2n-3\)~个 \(0\).
    按列~\(n\) 展开,
    并注意, 若 \(i < j\),
    则 \(U(n|{2n-1})\) 与 \(U({2n-1}|{2n-1})\)
    的 \((i, j)\)-元为 \(0\),
    我们得
    \begin{align*}
        \det {(U)}
        = {} &
        \hphantom{{} + {}}
        (-1)^{n+(2n-1)}\, [U]_{n,2n-1} \det {(U(n|{2n-1}))}
        \\
        &
        +
        (-1)^{(2n-1)+(2n-1)}\,
        [U]_{2n-1,2n-1} \det {(U({2n-1}|{2n-1}))}
        \\
        = {} &
        (-1)^{n-1} n a_0\, 1\, a_0^{n-2}\, (-n a_n)^{n-1}
        + a_{n-1}\, 1\, a_0^{n-2}\, ((1-n) a_{n-1})^{n-1}
        \\
        = {} &
        a_0^{n-2}
        (n^n a_0 a_n^{n-1} + (1-n)^{n-1} a_{n-1}^n).
    \end{align*}
    则 \(f\) 的判别式
    \begin{align*}
        (-1)^{n(n-1)/2} \det {(T(f; n))}
        = {} &
        (-1)^{n(n-1)/2} \det {(U)}
        \\
        = {} &
        (-1)^{n(n-1)/2} a_0^{n-2}
        (n^n a_0 a_n^{n-1} + (1-n)^{n-1} a_{n-1}^n).
    \end{align*}
    注意, 若 \(n = 2\), 则这与上个例的结果是一样的.
\end{example}

以下, 我们设假定~\ref{standing-assumption-of-resultant} 成立,
且 \(a_0 \neq 0\), \(b_0 \neq 0\),
\(m \geq 1\), 且 \(n \geq 1\).
设
\begin{align*}
    f(x) & = a_0 (x - x_1) (x - x_2) \dots (x - x_n), \\
    g(x) & = b_0 (x - y_1) (x - y_2) \dots (x - y_m).
\end{align*}
设 \(z\) 是参数.
设 \(r\) 是数.
考虑整式
\begin{align*}
    h(x; z, r)
    = {} &
    b_0 (z - r x - y_1) (z - r x - y_2) \dots
    (z - r x - y_m)
    \\
    = {} &
    \sum_{k = 0}^{m} {b_k (z - r x)^{m-k}}
    =
    \sum_{k = 0}^{m} {c_k (z, r) x^{m-k}},
\end{align*}
其中, \(c_k (z, r)\) 是关于
\(z\), \(r\), \(b_0\), \(b_1\), \(\dots\), \(b_m\)
的整式.
约定, 若 \(k < 0\) 或 \(k > m\),
则整式 \(c_k (z, r) = 0\).
作 \(f\) 与 \(h\) 的 Sylvester 阵 \(S(f, h; n, m)\).
它是 \(m + n\)~级阵, 且
\begin{align*}
    [S(f, h; n, m)]_{i,j} =
    \begin{cases}
        a_{i-j},            & j \leq m; \\
        c_{i-(j-m)} (z, r), & j > m.
    \end{cases}
\end{align*}
由结式的性质~(4), \(f\) 与 \(h\) 的结式
\begin{align*}
    \det {(S(f, h; n, m))}
    = {} &
    a_0^m h(x_1; z, r) h(x_2; z, r) \dots h(x_n; z, r)
    \\
    = {} &
    a_0^m b_0^n
    \prod_{i = 1}^{n} \prod_{j = 1}^{m}
    {(z - r x_i - y_j)}
    \\
    = {} &
    a_0^m b_0^n
    \prod_{i = 1}^{n} \prod_{j = 1}^{m}
    {(z - (y_j + r x_i))}.
\end{align*}
则这是一个关于 \(z\) 的 \(m n\)~次整式,
且它的 \(z^{m n}\) 项系数是 \(a_0^m b_0^n\).
不难看出, 它的全部的根是 \(y_j + r x_i\),
\(i = 1\), \(\dots\), \(n\),
且 \(j = 1\), \(\dots\), \(m\).

另考虑整式
\begin{align*}
    p(x; z) = b_0 (z - x y_1) (z - x y_2) \dots (z - x y_m).
\end{align*}
若 \(x \neq 0\), 则
\begin{align*}
    p(x; z)
    = {} &
    x^m b_0
    \left(\frac{z}{x} - y_1\right)
    \left(\frac{z}{x} - y_2\right)
    \dots
    \left(\frac{z}{x} - y_m\right)
    \\
    = {} &
    x^m \sum_{k = 0}^{m} {b_k \left(\frac{z}{x}\right)^{m-k}}
    \\
    = {} &
    x^m \sum_{k = 0}^{m} {b_k z^{m-k} x^{k-m}}
    \\
    = {} &
    \sum_{k = 0}^{m} {x^m b_k z^{m-k} x^{k-m}}
    \\
    = {} &
    \sum_{k = 0}^{m} {(b_k z^{m-k}) x^k}
    \\
    = {} &
    \sum_{\ell = 0}^{m} {(b_{m-\ell} z^\ell) x^{m-\ell}}.
\end{align*}
若 \(x = 0\), 则
\begin{align*}
    p(x; z)
    = b_0 (z - 0 y_1) (z - 0 y_2) \dots (z - 0 y_m)
    = b_0 z^m,
\end{align*}
且
\begin{align*}
    \sum_{\ell = 0}^{m} {(b_{m-\ell} z^\ell) x^{m-\ell}}
    = {} &
    \sum_{\ell = 0}^{m-1} {(b_{m-\ell} z^\ell) x^{m-\ell}}
    + (b_0 z^m) x^0
    \\
    = {} &
    \sum_{\ell = 0}^{m-1} {(b_{m-\ell} z^\ell) 0^{m-\ell}}
    + b_0 z^m
    \\
    = {} &
    b_0 z^m.
\end{align*}
所以, 对任何数 \(x\),
\begin{align*}
    p(x; z)
    = {} &
    b_0 (z - x y_1) (z - x y_2) \dots (z - x y_m)
    \\
    = {} &
    \sum_{\ell = 0}^{m} {(b_{m-\ell} z^\ell) x^{m-\ell}}
    =
    \sum_{\ell = 0}^{m} {d_\ell (z) x^{m-\ell}},
\end{align*}
其中,
\begin{align*}
    d_\ell (z) =
    \begin{cases}
        b_{m-\ell} z^\ell, & 0 \leq \ell \leq m; \\
        0,                 & \text{别的情形}.
    \end{cases}
\end{align*}
作 \(f\) 与 \(p\) 的 Sylvester 阵 \(S(f, p; n, m)\).
它是 \(m + n\)~级阵, 且
\begin{align*}
    [S(f, p; n, m)]_{i,j} =
    \begin{cases}
        a_{i-j},         & j \leq m; \\
        d_{i-(j-m)} (z), & j > m.
    \end{cases}
\end{align*}
由结式的性质~(4), \(f\) 与 \(p\) 的结式
\begin{align*}
    \det {(S(f, p; n, m))}
    = {} &
    a_0^m p(x_1; z) p(x_2; z) \dots p(x_n; z)
    \\
    = {} &
    a_0^m b_0^n
    \prod_{i = 1}^{n} \prod_{j = 1}^{m}
    {(z - x_i y_j)}.
\end{align*}
则这是一个关于 \(z\) 的 \(m n\)~次整式,
且它的 \(z^{m n}\) 项系数是 \(a_0^m b_0^n\).
不难看出, 它的全部的根是 \(x_i y_j\),
\(i = 1\), \(\dots\), \(n\),
且 \(j = 1\), \(\dots\), \(m\).

另考虑整式
\begin{align*}
    q(x; z)
    = {} &
    b_0 (x z - y_1) (x z - y_2) \dots (x z - y_m)
    \\
    = {} &
    \sum_{k = 0}^{m} b_k (x z)^{m-k}
    \\
    = {} &
    \sum_{k = 0}^{m} (b_k z^{m-k}) x^{m-k}
    \\
    = {} &
    \sum_{k = 0}^{m} {e_k (z) x^{m-k}},
\end{align*}
其中,
\begin{align*}
    e_k (z) =
    \begin{cases}
        b_k z^{m-k}, & 0 \leq k \leq m, \\
        0,           & \text{别的情形}.
    \end{cases}
\end{align*}
作 \(f\) 与 \(q\) 的 Sylvester 阵 \(S(f, q; n, m)\).
它是 \(m + n\)~级阵, 且
\begin{align*}
    [S(f, q; n, m)]_{i,j} =
    \begin{cases}
        a_{i-j},         & j \leq m; \\
        e_{i-(j-m)} (z), & j > m.
    \end{cases}
\end{align*}
由结式的性质~(4), \(f\) 与 \(q\) 的结式
\begin{align*}
    \det {(S(f, q; n, m))}
    = {} &
    a_0^m q(x_1; z) q(x_2; z) \dots q(x_n; z)
    \\
    = {} &
    a_0^m b_0^n
    \prod_{i = 1}^{n}
    \prod_{j = 1}^{m} {(x_i z - y_j)}.
\end{align*}
设 \(x_i \neq 0\), \(i = 1\), \(2\), \(\dots\), \(n\).
则
\begin{align*}
    \det {(S(f, q; n, m))}
    = {} &
    a_0^m b_0^n
    \prod_{i = 1}^{n}
    \prod_{j = 1}^{m} {\left( x_i \left( z - \frac{y_j}{x_i} \right) \right)}
    \\
    = {} &
    a_0^m b_0^n
    \prod_{i = 1}^{n} x_i^m
    \prod_{j = 1}^{m} {\left( z - \frac{y_j}{x_i} \right)}
    \\
    = {} &
    b_0^n a_0^m
    \prod_{i = 1}^{n} x_i^m
    \cdot
    \prod_{i = 1}^{n} \prod_{j = 1}^{m}
    {\left( z - \frac{y_j}{x_i} \right)}
    \\
    = {} &
    b_0^n a_0^m
    \left(\prod_{i = 1}^{n} x_i\right)^m
    \prod_{i = 1}^{n} \prod_{j = 1}^{m}
    {\left( z - \frac{y_j}{x_i} \right)}
    \\
    = {} &
    b_0^n
    \left( a_0 \prod_{i = 1}^{n} x_i \right)^m
    \prod_{i = 1}^{n} \prod_{j = 1}^{m}
    {\left( z - \frac{y_j}{x_i} \right)}.
\end{align*}
注意,
\begin{align*}
    a_0 x_1 x_2 \dots x_n
    = {} &
    (-1)^n a_0 (0 - x_1) (0 - x_2) \dots (0 - x_n)
    \\
    = {} &
    (-1)^n f(0)
    \\
    = {} &
    (-1)^n a_n.
\end{align*}
则
\begin{align*}
    \det {(S(f, q; n, m))}
    = (-1)^{m n} a_n^m b_0^n
    \prod_{i = 1}^{n} \prod_{j = 1}^{m}
    {\left( z - \frac{y_j}{x_i} \right)}.
\end{align*}
这是一个关于 \(z\) 的 \(m n\)~次整式,
且它的 \(z^{m n}\) 项系数是 \((-1)^{m n} a_n^m b_0^n\).
不难看出, 它的全部的根是 \(y_j / x_i\),
\(i = 1\), \(\dots\), \(n\),
且 \(j = 1\), \(\dots\), \(m\).

由前面的计算与讨论, 下面的结论是显然的.

\begin{theorem}
    设 \(n\), \(m\) 是正整数.
    设 \(a_0 \neq 0\), 且 \(b_0 \neq 0\).
    设
    \begin{align*}
        f(x) &
        = \sum_{k = 0}^{n} {a_k x^{n-k}}
        = a_0 (x - x_1) (x - x_2) \dots (x - x_n),
        \\
        g(x) &
        = \sum_{k = 0}^{m} {b_k x^{m-k}}
        = b_0 (x - y_1) (x - y_2) \dots (x - y_m).
    \end{align*}
    为方便, 我们约定:
    若 \(k < 0\) 或 \(k > n\), 则 \(a_k = 0\).
    设 \(z\) 是参数.

    (1)
    设 \(r\) 是数.
    作整式
    \begin{align*}
        h(x; z, r)
        = g(z - r x)
        = \sum_{k = 0}^{m} {b_k (z - r x)^{m-k}}
        = \sum_{k = 0}^{m} {c_k (z, r) x^{m-k}},
    \end{align*}
    其中, \(c_k (z, r)\) 是关于
    \(z\), \(r\), \(b_0\), \(b_1\), \(\dots\), \(b_m\)
    的整式,
    且若 \(k < 0\) 或 \(k > m\),
    则整式 \(c_k (z, r) = 0\).
    作 \(f\) 与 \(h\) 的 Sylvester 阵 \(S(f, h; n, m)\).
    它是 \(m + n\)~级阵, 且
    \begin{align*}
        [S(f, h; n, m)]_{i,j} =
        \begin{cases}
            a_{i-j},            & j \leq m; \\
            c_{i-(j-m)} (z, r), & j > m.
        \end{cases}
    \end{align*}
    则 \(f\) 与 \(h\) 的结式
    \begin{align*}
        \det {(S(f, h; n, m))}
        = a_0^m b_0^n
        \prod_{i = 1}^{n} \prod_{j = 1}^{m}
        {(z - (y_j + r x_i))}
    \end{align*}
    是一个关于 \(z\) 的 \(m n\)~次整式,
    且它的 \(z^{m n}\) 项系数是 \(a_0^m b_0^n\),
    且它的全部的根是 \(y_j + r x_i\),
    \(i = 1\), \(\dots\), \(n\),
    且 \(j = 1\), \(\dots\), \(m\).

    (2)
    作整式
    \begin{align*}
        p(x; z)
        = x^m g\left(\frac{z}{x}\right)
        = \sum_{\ell = 0}^{m} {(b_{m-\ell} z^\ell) x^{m-\ell}}
        = \sum_{\ell = 0}^{m} {d_\ell (z) x^{m-\ell}},
    \end{align*}
    其中, \(d_\ell (z)\) 是关于
    \(z\), \(b_0\), \(b_1\), \(\dots\), \(b_m\)
    的整式,
    且若 \(\ell < 0\) 或 \(\ell > m\),
    则整式 \(d_\ell (z, r) = 0\).
    作 \(f\) 与 \(p\) 的 Sylvester 阵 \(S(f, p; n, m)\).
    它是 \(m + n\)~级阵, 且
    \begin{align*}
        [S(f, p; n, m)]_{i,j} =
        \begin{cases}
            a_{i-j},         & j \leq m; \\
            d_{i-(j-m)} (z), & j > m.
        \end{cases}
    \end{align*}
    则 \(f\) 与 \(p\) 的结式
    \begin{align*}
        \det {(S(f, p; n, m))}
        = a_0^m b_0^n
        \prod_{i = 1}^{n} \prod_{j = 1}^{m}
        {(z - x_i y_j)}
    \end{align*}
    是一个关于 \(z\) 的 \(m n\)~次整式,
    且它的 \(z^{m n}\) 项系数是 \(a_0^m b_0^n\),
    且它的全部的根是 \(x_i y_j\),
    \(i = 1\), \(\dots\), \(n\),
    且 \(j = 1\), \(\dots\), \(m\).

    (3)
    设 \(x_i \neq 0\), \(i = 1\), \(2\), \(\dots\), \(n\).
    作整式
    \begin{align*}
        q(x; z)
        = g (z x)
        = \sum_{k = 0}^{m} (b_k z^{m-k}) x^{m-k}
        = \sum_{k = 0}^{m} {e_k (z) x^{m-k}},
    \end{align*}
    其中, \(e_k (z)\) 是关于
    \(z\), \(b_0\), \(b_1\), \(\dots\), \(b_m\)
    的整式,
    且若 \(k < 0\) 或 \(k > m\),
    则整式 \(e_k (z, r) = 0\).
    作 \(f\) 与 \(q\) 的 Sylvester 阵 \(S(f, q; n, m)\).
    它是 \(m + n\)~级阵, 且
    \begin{align*}
        [S(f, q; n, m)]_{i,j} =
        \begin{cases}
            a_{i-j},         & j \leq m; \\
            e_{i-(j-m)} (z), & j > m.
        \end{cases}
    \end{align*}
    则 \(f\) 与 \(q\) 的结式
    \begin{align*}
        \det {(S(f, q; n, m))}
        = (-1)^{m n} a_n^m b_0^n
        \prod_{i = 1}^{n} \prod_{j = 1}^{m}
        {\left( z - \frac{y_j}{x_i} \right)}
    \end{align*}
    是一个关于 \(z\) 的 \(m n\)~次整式,
    且它的 \(z^{m n}\) 项系数是 \((-1)^{m n} a_n^m b_0^n\),
    且它的全部的根是 \(y_j / x_i\),
    \(i = 1\), \(\dots\), \(n\),
    且 \(j = 1\), \(\dots\), \(m\).
\end{theorem}

\begin{example}
    设
    \begin{align*}
        f(x) & = x^2 - x - 1 = (x - x_1) (x - x_2), \\
        g(x) & = x^2 + x + 1 = (x - y_1) (y - y_2).
    \end{align*}
    我们求以
    \(y_j - x_i\), \(i\), \(j = 1\), \(2\)
    为根的整式,
    与以
    \(x_i y_j\), \(i\), \(j = 1\), \(2\)
    为根的整式.
    注意,
    \begin{align*}
        h(x; z, -1)
        = {} &
        g(z - (-1) x)
        \\
        = {} &
        (z + x)^2 + (z + x) + 1
        \\
        = {} &
        (z^2 + 2 z x + x^2) + (z + x) + 1
        \\
        = {} &
        1 x^2 + (2z + 1) x + (z^2 + z + 1),
    \end{align*}
    且
    \begin{align*}
        p(x; z)
        = x^2 g\left( \frac{z}{x} \right)
        = x^2 \left(\frac{z^2}{x^2} + \frac{z}{x} + 1\right)
        = 1 x^2 + z x + z^2.
    \end{align*}
    由此, 我们写出 \(f\), \(h\) 的 Sylvester 阵
    \begin{align*}
        S(f, h; 2, 2) =
        \begin{bmatrix}
            1  & 0  & 1       & 0       \\
            -1 & 1  & 2z+1    & 1       \\
            -1 & -1 & z^2+z+1 & 2z+1    \\
            0  & -1 & 0       & z^2+z+1 \\
        \end{bmatrix},
    \end{align*}
    与 \(f\), \(p\) 的 Sylvester 阵
    \begin{align*}
        S(f, p; 2, 2) =
        \begin{bmatrix}
            1  & 0  & 1   & 0   \\
            -1 & 1  & z   & 1   \\
            -1 & -1 & z^2 & z   \\
            0  & -1 & 0   & z^2 \\
        \end{bmatrix}.
    \end{align*}
    注意, 这 \(2\)~个 Sylvester 阵的元都%
    全是关于 \(z\) 的, 系数全为整数的整式,
    所以它们的行列式也都是%
    关于 \(z\) 的, 系数全为整数的整式.

    注意,
    \begin{align*}
        \det {
            \begin{bmatrix}
                1  & 0  & a & 0 \\
                -1 & 1  & b & a \\
                -1 & -1 & c & b \\
                0  & -1 & 0 & c \\
            \end{bmatrix}
        }
        = {} &
        \det {
            \begin{bmatrix}
                1    & 0  & a   & 0 \\
                -1+1 & 1  & b+a & a \\
                -1+1 & -1 & c+a & b \\
                0    & -1 & 0   & c \\
            \end{bmatrix}
        }
        \\
        = {} &
        \det {
            \begin{bmatrix}
                1 & 0  & a   & 0 \\
                0 & 1  & b+a & a \\
                0 & -1 & c+a & b \\
                0 & -1 & 0   & c \\
            \end{bmatrix}
        }
        \\
        = {} &
        \det {
            \begin{bmatrix}
                1 & 0  & a           & 0   \\
                0 & 1  & b+a         & a   \\
                0 & -1+1 & c+a+(b+a) & b+a \\
                0 & -1+1 & 0+(b+a)   & c+a \\
            \end{bmatrix}
        }
        \\
        = {} &
        \det {
            \begin{bmatrix}
                1 & 0 & a      & 0   \\
                0 & 1 & a+b    & a   \\
                0 & 0 & 2a+b+c & a+b \\
                0 & 0 & a+b    & a+c \\
            \end{bmatrix}
        }
        \\
        = {} &
        1
        \det {
            \begin{bmatrix}
                1 & a+b    & a   \\
                0 & 2a+b+c & a+b \\
                0 & a+b    & a+c \\
            \end{bmatrix}
        }
        \\
        = {} &
        1 \cdot 1
        \det {
            \begin{bmatrix}
                2a+b+c & a+b \\
                a+b    & a+c \\
            \end{bmatrix}
        }
        \\
        = {} &
        \det {
            \begin{bmatrix}
                2a+b+c-(a+b) & a+b-(a+c) \\
                a+b          & a+c       \\
            \end{bmatrix}
        }
        \\
        = {} &
        \det {
            \begin{bmatrix}
                a+c & -(c-b) \\
                a+b & a+c    \\
            \end{bmatrix}
        }
        \\
        = {} &
        (c + a)^2 + (b + a) (c - b).
    \end{align*}

    由前面的计算, 与上个定理,
    以
    \(y_j - x_i\), \(i\), \(j = 1\), \(2\)
    为根的整式是
    \begin{align*}
        \det {(S(f, h; 2, 2))}
        = {} &
        (z^2 + z + 1 + 1)^2
        + (2 z + 1 + 1) (z^2 + z + 1 - (2 z + 1))
        \\
        = {} &
        z^4 + 4 z^3 + 5 z^2 + 2 z + 4,
    \end{align*}
    且以
    \(x_i y_j\), \(i\), \(j = 1\), \(2\)
    为根的整式是
    \begin{align*}
        \det {(S(f, p; 2, 2))}
        = {} &
        (z^2 + 1)^2 + (z + 1) (z^2 - z)
        \\
        = {} &
        z^4 + z^3 + 2 z^2 - z + 1.
    \end{align*}

    我们或许能体会到上个定理的好的地方:
    因为 \(f\), \(h\), \(p\) 是关于 \(x\), \(z\) 的,
    系数全为整数的整式,
    即使不具体地计算, 我们也知道,
    存在关于 \(z\) 的, 系数全为整数的整式,
    其的根是 \(g\), \(f\) 的根的差 (或积).
    % 最后, 我直接地计算
    % \begin{align*}
    %     &
    %     (z - (y_1 - x_1)) (z - (y_1 - x_2))
    %     (z - (y_2 - x_1)) (z - (y_2 - x_2))
    %     \\
    %     = {} &
    %     (z - y_1 + x_1) (z - y_1 + x_2)
    %     (z - y_2 + x_1) (z - y_2 + x_2)
    %     \\
    %     = {} &
    %     ((z - y_1)^2 + (z - y_1) (x_1 + x_2) + x_1 x_2)
    %     ((z - y_2)^2 + (z - y_2) (x_1 + x_2) + x_1 x_2)
    %     \\
    %     = {} &
    %     ((z - y_1)^2 + (z - y_1) - 1)
    %     ((z - y_2)^2 + (z - y_2) - 1)
    %     \\
    %     = {} &
    %     (z^2 - 2 z y_1 + y_1^2 + z - y_1 - 1)
    %     (z^2 - 2 z y_2 + y_2^2 + z - y_2 - 1)
    %     \\
    %     = {} &
    %     (z^2 - 2 z y_1 + z - 2 y_1 - 2)
    %     (z^2 - 2 z y_2 + z - 2 y_2 - 2)
    %     \\
    %     = {} &
    %     ((z^2 + z - 2) - 2 y_1 (z + 1))
    %     ((z^2 + z - 2) - 2 y_2 (z + 1))
    %     \\
    %     = {} &
    %     (z^2 + z - 2)^2 - 2 (z^2 + z - 2) (z + 1) (y_1 + y_2)
    %     + 4 (z + 1)^2 y_1 y_2
    %     \\
    %     = {} &
    %     (z^2 + z - 2)^2 + 2 (z^2 + z - 2) (z + 1) + 4 (z + 1)^2
    %     \\
    %     = {} &
    %     z^4 + 4 z^3 + 5 z^2 + 2 z + 4,
    % \end{align*}
    % 与
    % \begin{align*}
    %     &
    %     (z - x_1 y_1) (z - x_1 y_2) (z - x_2 y_1) (z - x_2 y_2)
    %     \\
    %     = {} &
    %     (z^2 - z x_1 (y_1 + y_2) + x_1^2 y_1 y_2)
    %     (z^2 - z x_2 (y_1 + y_2) + x_2^2 y_1 y_2)
    %     \\
    %     = {} &
    %     (z^2 + z x_1 + x_1^2)
    %     (z^2 + z x_2 + x_2^2)
    %     \\
    %     = {} &
    %     (z^2 + z x_1 + 1 + x_1)
    %     (z^2 + z x_2 + 1 + x_2)
    %     \\
    %     = {} &
    %     ((z^2 + 1) + (z + 1) x_1)
    %     ((z^2 + 1) + (z + 1) x_2)
    %     \\
    %     = {} &
    %     (z^2 + 1)^2 + (z^2 + 1) (z + 1) (x_1 + x_2) + (z + 1)^2 x_1 x_2
    %     \\
    %     = {} &
    %     (z^2 + 1)^2 + (z^2 + 1) (z + 1) - (z + 1)^2
    %     \\
    %     = {} &
    %     z^4 + z^3 + 2 z^2 - z + 1,
    % \end{align*}
    % 其中, 我用了
    % \begin{align*}
    %     x_1 + x_2 = -\frac{-1}{1} = 1,
    %     & \quad
    %     x_1 x_2 = \frac{-1}{1} = -1,
    %     \\
    %     y_1 + y_2 = -\frac{1}{1} = -1,
    %     & \quad
    %     y_1 y_2 = \frac{1}{1} = 1;
    % \end{align*}
    % 注意, 若
    % \(a x^2 + b x + c = a (x - x_1) (x - x_2)\),
    % 则
    % \begin{align*}
    %     a x^2 + b x + c
    %     = {} &
    %     a (x - x_1) (x - x_2)
    %     \\
    %     = {} &
    %     a (x^2 - (x_1 + x_2) x + x_1 x_2)
    %     \\
    %     = {} &
    %     a x^2 + (-a (x_1 + x_2)) x + a x_1 x_2,
    % \end{align*}
    % 故
    % \begin{align*}
    %     x_1 + x_2 = \frac{-a (x_1 + x_2)}{-a}
    %     = \frac{b}{-a} = -\frac{b}{a},
    %     \quad
    %     x_1 x_2 = \frac{a x_1 x_2}{a} = \frac{c}{a}.
    % \end{align*}
    % 由此, 您可以看到, 这验证了上个定理.
\end{example}

\begin{theorem}
    设 \(n\)~是正整数.
    设 \(a_0 \neq 0\).
    设
    \begin{align*}
        f(x) &
        = \sum_{k = 0}^{n} {a_k x^{n-k}}
        = a_0 (x - x_1) (x - x_2) \dots (x - x_n).
    \end{align*}
    为方便, 我们约定:
    若 \(k < 0\) 或 \(k > n\), 则 \(a_k = 0\).
    设 \(z\) 是参数.
    设 \(u(x)\), \(v(x)\) 是整式,
    且 \(u(x_i) \neq 0\), \(i = 1\), \(\dots\), \(n\).
    设
    \begin{align*}
        r(x; z) = u(x) z - v(x)
        = \sum_{k = 0}^{m} {b_k (z) x^{m-k}},
    \end{align*}
    其中, \(b_k (z)\) 是关于 \(z\) 的整式,
    且若 \(k < 0\) 或 \(k > m\),
    则整式 \(b_k (z) = 0\).
    作 \(f\) 与 \(r\) 的 Sylvester 阵 \(S(f, r; n, m)\).
    它是 \(m + n\)~级阵, 且
    \begin{align*}
        [S(f, r; n, m)]_{i,j} =
        \begin{cases}
            a_{i-j},         & j \leq m; \\
            b_{i-(j-m)} (z), & j > m.
        \end{cases}
    \end{align*}
    则 \(f\) 与 \(r\) 的结式
    \begin{align*}
        \det {(S(f, r; n, m))}
        = a_0^m \left(\prod_{i = 1}^{n} {u(x_i)} \right)
        \prod_{i = 1}^{n} {\left(z - \frac{v(x_i)}{u(x_i)}\right)}
    \end{align*}
    是一个关于 \(z\) 的 \(n\)~次整式,
    且它的 \(z^{n}\) 项系数是 \(a_0^m u(x_1) u(x_2) \dots u(x_n)\),
    且它的全部的根是 \(v(x_i) / u(x_i)\),
    \(i = 1\), \(\dots\), \(n\).
\end{theorem}

\begin{proof}
    我们计算结式即可;
    由此我们得别的结论.

    由结式的性质~(4), \(f\) 与 \(r\) 的结式
    \begin{align*}
        \det {(S(f, r; n, m))}
        = {} &
        a_0^m \prod_{i = 1}^{n} {r(x_i; z)}
        \\
        = {} &
        a_0^m \prod_{i = 1}^{n} {(u(x_i) z - v(x_i))}
        \\
        = {} &
        a_0^m \prod_{i = 1}^{n}
        {\left(u(x_i) \left(z - \frac{v(x_i)}{u(x_i)}\right) \right)}
        \\
        = {} &
        a_0^m \left(\prod_{i = 1}^{n} {u(x_i)} \right)
        \prod_{i = 1}^{n} {\left(z - \frac{v(x_i)}{u(x_i)}\right)}.
        \qedhere
    \end{align*}
\end{proof}

值得提, 前面, 在研究问题
``若数 \(c\) 适合 \(f(c) = 0\),
则 \(y = v(c) / u(c)\) 是何''
(其中, \(f(x)\), \(u(x)\), \(v(x)\) 是整式,
且 \(u(c) \neq 0\))
时,
我们其实用了一样的结式
(即 \(f(x)\) 与 \(u(x) y - v(x)\) 的结式).
由此可见, 结式的应用是多的与广的.

\end{document}

This material discusses
the resultant and the matrix of Sylvester.
